% archon:covers AlgebraicJacobian/Albanese/CodimOneExtension.lean
% NOTE: Weil-divisor API extension of RiemannRoch_WeilDivisor.tex --- re-uses
% Scheme.WeilDivisor, Scheme.PrimeDivisor, Scheme.RationalMap.order from
% \cref{chap:RiemannRoch_WeilDivisor}. The order map
% \texttt{AlgebraicGeometry.Scheme.RationalMap.order} is not re-defined here;
% see \cref{def:order_at_point} for the canonical reference.

\chapter{Codimension-1 indeterminacy extension (A.4.a)}
\label{chap:Albanese_CodimOneExtension}

% STRATEGY NOTE (iter-174). This chapter is the A.4.a sub-row of Route A
% (the positive-genus arm of \texttt{nonempty\_jacobianWitness}). Per
% STRATEGY.md L29 it is the \emph{risk-dominant} A.4 sub-row
% ("project-fatal if it stalls"). The chapter has two outputs:
% (i) \cref{thm:codim_one_extension} --- Milne Theorem~3.1: a rational map
%     from a normal variety to a complete variety is defined on the complement
%     of a codimension-\(\geq 2\) closed subset. Specialised below to the
%     smooth-variety setting where the prover can phrase it as
%     "no codim-1 indeterminacy \(\Rightarrow\) extension". This is one of the
%     two inputs to Milne Theorem~3.2 in the sibling chapter
%     \cref{chap:Albanese_Thm32RationalMapExtension}.
% (ii) \cref{lem:milne_codim1_indeterminacy} --- Milne Lemma~3.3: for a
%     rational map from a nonsingular variety to a group variety, the
%     indeterminacy locus is either empty or of pure codimension \(1\). This is
%     the second input to Milne Theorem~3.2; it consumes nothing from A.4.b.
% A bonus output is \cref{thm:weil_divisor_obstruction}: the Weil-divisor
% reformulation of (i) that a smooth-surface codim-2 extension is obstructed
% exactly by the order of the pulled-back coordinates of \(f\) being negative
% at a prime divisor. This is the lemma A.4.a uses to align with the project's
% \cref{chap:RiemannRoch_WeilDivisor} machinery.
%
% Depth-\(\geq 2\) extension on a smooth surface, used in the proof body of
% \cref{thm:codim_one_extension} when reducing the general-codim case to the
% normal-DVR-codim-1 case, consumes
% \cref{cor:regular_cohen_macaulay} of
% \cref{chap:Albanese_AuslanderBuchsbaum} (the A.4.b sub-row). The chapter
% records the dependency but does not re-prove it.
%
% The Lean target file \texttt{AlgebraicJacobian/Albanese/CodimOneExtension.lean}
% does not yet exist; this chapter is the prover-ready specification for
% creating it.

\section{Setup and motivation}
\label{sec:codim1_setup}

Fix an algebraically closed field \(\bar k\) throughout. By a \emph{variety} over \(\bar k\)
we mean a geometrically integral, separated scheme of finite type over \(\bar k\)
(Milne's convention, \emph{Abelian Varieties}, p.~v); a variety is \emph{nonsingular}
(or \emph{smooth}) if its local ring at every closed point is regular. A variety is
\emph{complete} if it is proper over \(\bar k\).

A \emph{rational map} \(f \colon X \dashrightarrow Y\) between varieties is the equivalence
class of a regular morphism \(\varphi_U \colon U \to Y\) defined on a dense open subset
\(U \subseteq X\), where two pairs \((U, \varphi_U)\), \((U', \varphi_{U'})\) are equivalent
if \(\varphi_U\) and \(\varphi_{U'}\) agree on \(U \cap U'\) (Milne~\S I.3, p.~16). The map \(f\)
is \emph{defined at} a point \(x \in X\) if there is some representative \((U, \varphi_U)\)
with \(x \in U\); the set of such \(x\) is an open subset \(\mathrm{Dom}(f) \subseteq X\),
the \emph{domain of definition} of \(f\).

The motivating question of this chapter, following Milne~\S I.3, is:
\begin{quote}
  \emph{When does a rational map \(f \colon X \dashrightarrow Y\) extend to a regular
  morphism on all of \(X\)?}
\end{quote}
Without hypotheses on \(X\) or \(Y\) the question has no positive answer (Milne lists three
counterexamples on p.~16: the inclusion of a proper open into the ambient variety; the
rational map to the cuspidal cubic whose inverse fails to extend; the inverse of a
blow-up at a smooth point). The two structural inputs that buy unconditional extension
in our setting are:
\begin{itemize}
  \item \emph{Target completeness:} if \(Y\) is complete (proper over \(\bar k\)), the
    valuative criterion of properness forces the rational map to extend across every
    codim-\(1\) point of a normal source \(X\) (Milne Theorem~3.1).
  \item \emph{Target group structure:} if \(Y\) is a group variety, the difference map
    \(\Phi(x, y) := f(x) \cdot f(y)^{-1}\) on \(X \times X\) has a vanishing/pole divisor
    along the diagonal that controls the codimension of the indeterminacy locus
    (Milne Lemma~3.3).
\end{itemize}
For \(Y\) an abelian variety both conditions hold, and combining them yields Milne's
Theorem~3.2 (the headline of \cref{chap:Albanese_Thm32RationalMapExtension}). The
present chapter packages the two inputs separately:
\cref{thm:codim_one_extension} is Milne~3.1, and
\cref{lem:milne_codim1_indeterminacy} is Milne~3.3. The Weil-divisor reformulation
\cref{thm:weil_divisor_obstruction} translates the codim-\(1\) extension obstruction
in Milne~3.1's proof into the order-at-a-prime-divisor language pinned in
\cref{chap:RiemannRoch_WeilDivisor}, which is the form the Lean prover consumes when
working with the project's \(\mathrm{Scheme.WeilDivisor}\) API.

\section{The indeterminacy locus of a rational map}
\label{sec:indeterm_locus}

\begin{definition}


\leanok
  [Indeterminacy locus of a rational map]
  \label{def:indeterminacy_locus}
  \lean{AlgebraicGeometry.Scheme.RationalMap.indeterminacyLocus}
  % SOURCE: [Milne, Abelian Varieties], \S I.3, p.~16, the paragraph "Let \(V\) and \(W\)
  % be varieties over \(k\)\ldots" defining the domain of definition of a rational map
  % (read from references/abelian-varieties.pdf, PDF page 22).
  % SOURCE QUOTE: "A rational map \(\varphi\) is said to be \emph{defined} at a point
  % \(v\) of \(V\) if \(v \in U\) for some \((U, \varphi_U) \in \varphi\). The set \(U_1\) of
  % \(v\) at which \(\varphi\) is defined is open, and there is a regular map
  % \(\varphi_1 \colon U_1 \to W\) such that \((U_1, \varphi_1) \in \varphi\) --- clearly,
  % \(U_1 = \bigcup_{(U, \varphi_U) \in \varphi} U\) and we can define \(\varphi_1\) to be
  % the regular map such that \(\varphi_1 | U = \varphi_U\) for all $(U, \varphi_U) \in
  % \varphi$."
  \textit{Source: Milne, \emph{Abelian Varieties}, \S I.3, p.~16.}
  Let \(f \colon X \dashrightarrow Y\) be a rational map between varieties over \(\bar k\).
  The \emph{domain of definition} of \(f\) is the open set
  \[
    \mathrm{Dom}(f) \;:=\; \bigcup_{(U, \varphi_U) \in f} U \;\subseteq\; X,
  \]
  i.e.\ the union of the dense opens of all representatives of \(f\). The
  \emph{indeterminacy locus} of \(f\) is the closed complement
  \[
    Z(f) \;:=\; X \smallsetminus \mathrm{Dom}(f) \;\subseteq\; X.
  \]
  The map \(f\) is regular (i.e.\ defined on all of \(X\)) if and only if \(Z(f) = \emptyset\).

  \paragraph{Lean encoding scope.} The Lean target
  \texttt{AlgebraicGeometry.Scheme.RationalMap.indeterminacyLocus} is a method on the
  Mathlib \texttt{Scheme.RationalMap} structure pinned at
  \texttt{Mathlib.AlgebraicGeometry.Birational.RationalMap}; it returns a
  \texttt{Set X} (the underlying carrier of the closed complement of
  \texttt{Scheme.RationalMap.domain}, which is already exposed by Mathlib). The
  closedness of \(Z(f)\) is a separate \texttt{IsClosed} witness, packaged either as a
  field of \texttt{indeterminacyLocus} or as a standalone lemma
  \texttt{Scheme.RationalMap.isClosed\_indeterminacyLocus}, whichever is cleaner with
  Mathlib's existing API at commit \texttt{b80f227}.
\end{definition}

\begin{lemma}\leanok
[Indeterminacy locus is closed]
  \label{lem:isClosed_indeterminacy_locus}
  \lean{AlgebraicGeometry.Scheme.RationalMap.isClosed_indeterminacyLocus}
  \uses{def:indeterminacy_locus}
  For any rational map \(f \colon X \dashrightarrow Y\) of schemes, the
  indeterminacy locus \(Z(f) = X \smallsetminus \mathrm{Dom}(f)\)
  (\cref{def:indeterminacy_locus}) is a closed subset of \(X\): it is the
  complement of the open domain of definition \(\mathrm{Dom}(f)\).
\end{lemma}

\begin{proof}
  Proved directly in Lean: \(\mathrm{Dom}(f)\) is open, so its complement is closed.
\end{proof}

\begin{definition}


\leanok
  [Codimension-1-indeterminacy-free rational maps]
  \label{def:codim_one_indeterminacy}
  \lean{AlgebraicGeometry.Scheme.RationalMap.CodimOneFree}
  \uses{def:indeterminacy_locus, def:prime_divisor}
  % SOURCE: [Milne, Abelian Varieties], Theorem~3.1, \S I.3, p.~16 --- the conclusion
  % "an open subset \(U\) of \(V\) whose complement \(V \smallsetminus U\) has codimension \(\geq 2\)"
  % is the standing predicate we abstract (read from references/abelian-varieties.pdf,
  % PDF page 22).
  % SOURCE QUOTE: "A rational map \(\varphi \colon V \dashrightarrow W\) from a normal
  % variety \(V\) to a complete variety \(W\) is defined on an open subset \(U\) of \(V\) whose
  % complement \(V \smallsetminus U\) has codimension \(\geq 2\)."
  \textit{Source: Milne, \emph{Abelian Varieties}, Theorem~3.1, \S I.3, p.~16 (the
  conclusion).}
  Let \(f \colon X \dashrightarrow Y\) be a rational map between varieties over \(\bar k\).
  We say \(f\) is \emph{codim-1-indeterminacy-free} (or has \emph{codimension-\(\geq 2\)
  indeterminacy}) if every irreducible component of the closed subset \(Z(f) \subseteq X\)
  has codimension \(\geq 2\) in \(X\). Equivalently, no prime divisor of \(X\) is contained in
  \(Z(f)\); equivalently, every point \(\eta \in X\) with \(\mathrm{codim}_X \overline{\{\eta\}} = 1\)
  lies in \(\mathrm{Dom}(f)\).

  \paragraph{Why the abstraction.} Milne's Theorem~3.1 (\cref{thm:codim_one_extension}
  below) and Lemma~3.3 (\cref{lem:milne_codim1_indeterminacy} below) both produce
  conclusions about the codimension of \(Z(f)\). Phrasing the codim-\(\geq 2\) condition
  as a named predicate keeps the cited statements concise and lets the consumer
  \cref{chap:Albanese_Thm32RationalMapExtension} read off the combination
  (\(Z(f)\) has codim \(\geq 2\) from \cref{thm:codim_one_extension} \emph{and} \(Z(f)\) is
  empty or pure codim 1 from \cref{lem:milne_codim1_indeterminacy} \(\Rightarrow\) \(Z(f)\)
  is empty) as a single line.

  \paragraph{Lean encoding scope.} The Lean target
  \texttt{AlgebraicGeometry.Scheme.RationalMap.CodimOneFree} is a \texttt{Prop}-valued
  predicate, encoded via Mathlib's \texttt{Order.coheight} on the specialisation
  preorder of \texttt{X.carrier} as in
  \cref{def:prime_divisor}: \(f\) is codim-1-indeterminacy-free if and only if for every
  \(\eta \in X\) with \(\texttt{Order.coheight\ }\eta = 1\), \(\eta \in
  \texttt{f.domain}\). The encoding matches the standing \(\mathrm{Order.coheight}\)
  convention pinned in \cref{def:prime_divisor}, so prime divisors of \(X\) and the
  hypothesis ``no codim-\(1\) indeterminacy'' use the same Mathlib idiom.
\end{definition}

\section{Smoothness yields a DVR at every codim-1 point}
\label{sec:smooth_dvr}

The codim-\(1\) extension argument in Milne's proof of Theorem~3.1 hinges on the local
ring at the generic point of a prime divisor being a discrete valuation ring (so that
the valuative criterion of properness applies). On a normal variety this is the
defining property of regularity in codimension one; on a smooth variety it is automatic.

The smooth-implies-regular-local-stalk implication itself is the Jacobian-criterion
direction of Stacks tag \texttt{00TT}; we package it as a separate sub-lemma so that
\cref{lem:smooth_codim_one_dvr} below reads cleanly as ``regular local of dimension
\(1\) is a DVR'', with the regularity prerequisite supplied by the sub-lemma.

\subsection{Stages 5a/5b: K\"ahler-differentials localisation substrate (iter-193)}
\label{subsec:kaehler_localisation_substrate}

The Jacobian-criterion direction of Stacks tag \texttt{00TT}
(\cref{lem:smooth_to_regular_local_ring} below) factors on the Lean side
through a six-stage chain. Stages~1--2 (smooth \(\Rightarrow\) flat at a
stalk; smooth \(\Rightarrow\) admits a standard-smooth presentation
locally) are axiom-clean since iter-191; Stages~3--4 (the
scheme-to-algebra bridge for standard-smooth presentations; freeness of
\(\Omega_{S/R}\) for a standard-smooth algebra \(R \to S\)) are
axiom-clean since iter-192. The remaining Stages~5--6 reduce the
Stacks~00TT conclusion to a freeness plus rank computation on K\"ahler
differentials, transported through localisation at the prime defining
the stalk. Iter-193 packaged Stages~5a (transport of freeness) and 5b
(matching rank under localisation) as the two standalone helpers below;
Stage~6 (the remaining Stacks~00OE smooth-algebra dimension formula
plus Stacks~02JK cotangent--K\"ahler-over-a-field bridge) is the open
Mathlib gap, decomposed iter-198 into the three named sub-lemmas
\cref{lem:smooth_algebra_krull_dim_formula} (6.A),
\cref{lem:cotangent_kahler_over_field} (6.B), and
\cref{lem:stage6_regular_stalk_assembly} (6.C) in
\cref{subsec:stage6_subgap_decomposition} immediately following.

Stages~1--4 of the chain are recorded below as named substrate lemmas, feeding
\cref{lem:smooth_to_regular_local_ring}.

\begin{lemma}\leanok
[Stage 1 --- smooth implies flat at every stalk]
  \label{lem:stalkMap_flat_of_smooth}
  \lean{AlgebraicGeometry.Scheme.stalkMap_flat_of_smooth}
  Let \(X \to \Spec(\bar k)\) be a smooth morphism over an algebraically closed
  field \(\bar k\). Then for every point \(z \in X\) the induced ring map on
  stalks is flat.
\end{lemma}

\begin{proof}
  Proved directly in Lean: the smooth-implies-flat instance composed with the
  stalk-map flatness characterisation. A sub-step of
  \cref{lem:smooth_to_regular_local_ring}.
\end{proof}

\begin{lemma}\leanok
[Stage 2 --- smooth implies a standard-smooth presentation at a point]
  \label{lem:exists_standardSmooth_at_of_smooth}
  \lean{AlgebraicGeometry.Scheme.exists_isStandardSmooth_at_of_smooth}
  For a smooth morphism \(X \to \Spec(\bar k)\) and a point \(z \in X\), there
  exist affine open neighbourhoods \(U \subseteq \Spec(\bar k)\) and \(V \ni z\)
  in \(X\) with \(V \subseteq X^{-1}(U)\) such that the induced ring map
  \(\Gamma(\Spec \bar k, U) \to \Gamma(X, V)\) is standard smooth (Stacks 00T7).
\end{lemma}

\begin{proof}
  Proved directly in Lean: re-export of the scheme-level Stacks 00T7 packaging.
  A sub-step of \cref{lem:exists_algebra_standardSmooth_stalk_localization}.
\end{proof}

\begin{lemma}\leanok
[Stage 3 --- standard-smooth algebra package at the stalk]
  \label{lem:exists_algebra_standardSmooth_stalk_localization}
  \lean{AlgebraicGeometry.Scheme.exists_algebra_isStandardSmooth_section_stalk_isLocalization_of_smooth}
  \uses{lem:exists_standardSmooth_at_of_smooth, lem:isLocalization_atPrime_stalk_of_affineOpen}
  For a smooth morphism \(X \to \Spec(\bar k)\) and a point \(z \in X\), there is
  an affine open neighbourhood \(V \ni z\) and an \(\Gamma(\Spec \bar k, U)\)-algebra
  structure on \(\Gamma(X, V)\) making it a standard-smooth algebra, together with a
  witness that the stalk \(\mathcal O_{X, z}\) is the localisation of \(\Gamma(X, V)\)
  at the prime ideal corresponding to \(z\).
\end{lemma}

\begin{proof}
  Proved directly in Lean: composes the Stage-2 presentation
  \cref{lem:exists_standardSmooth_at_of_smooth} with the ring-hom-to-algebra bridge
  and the affine-open stalk localisation
  \cref{lem:isLocalization_atPrime_stalk_of_affineOpen}. A sub-step of
  \cref{lem:smooth_to_regular_local_ring}.
\end{proof}

\begin{lemma}\leanok
[Stage 4 --- K\"ahler differentials of a standard-smooth algebra are free]
  \label{lem:module_free_kaehler_standardSmooth}
  \lean{AlgebraicGeometry.Scheme.module_free_kaehlerDifferential_of_isStandardSmooth}
  For a standard-smooth \(R\)-algebra \(S\), the module of K\"ahler differentials
  \(\Omega_{S/R}\) is free as an \(S\)-module (Stacks 00T7(2)), with basis the family
  \(\{ds_i\}\) of the standard-smooth presentation.
\end{lemma}

\begin{proof}
  Proved directly in Lean: from the basis provided by the standard-smooth
  presentation. Supplies the freeness hypothesis of
  \cref{lem:module_free_kaehler_localization}.
\end{proof}

\begin{lemma}


\leanok
  [Stage 5a --- K\"ahler-differentials freeness transports through localisation]
  \label{lem:module_free_kaehler_localization}
  \lean{AlgebraicGeometry.Scheme.module_free_kaehlerDifferential_localization}
  \uses{lem:module_free_kaehler_standardSmooth}
  Let \(R\) and \(S\) be commutative rings equipped with an \(R\)-algebra
  structure on \(S\), and assume that the module of K\"ahler differentials
  \(\Omega_{S/R}\) is free as an \(S\)-module. Let \(M \subseteq S\) be a
  submonoid and let \(S_M\) be a localisation of \(S\) at \(M\), so that
  \(R \to S \to S_M\) is a tower of \(R\)-algebras. Then the module of
  K\"ahler differentials \(\Omega_{S_M / R}\) is free as an
  \(S_M\)-module.
\end{lemma}

\begin{proof}
  The canonical comparison map
  \(\Omega_{S/R} \to \Omega_{S_M/R}\) induced by the localisation map
  \(S \to S_M\) realises \(\Omega_{S_M/R}\) as the localisation of the
  \(S\)-module \(\Omega_{S/R}\) at \(M\); i.e.\ it satisfies Mathlib's
  \texttt{IsLocalizedModule} predicate along the underlying ring
  localisation. This is Stacks tag \texttt{02JK}, packaged in Mathlib as
  \texttt{KaehlerDifferential.isLocalizedModule\_map}. Freeness of a
  module is preserved under localisation of modules via Mathlib's
  \texttt{Module.free\_of\_isLocalizedModule}; combining the two yields
  the conclusion.
\end{proof}

\begin{lemma}


\leanok
  [Stage 5b --- K\"ahler-differentials rank under localisation equals
  the relative dimension]
  \label{lem:rank_kaehler_localization_eq_relative_dim}
  \lean{AlgebraicGeometry.Scheme.rank_kaehlerDifferential_localization_eq_relativeDimension}
  \uses{lem:module_free_kaehler_localization}
  With notation as in \cref{lem:module_free_kaehler_localization},
  assume in addition that the \(R\)-algebra \(S\) is \emph{standard
  smooth of relative dimension \(n\)} (Mathlib's
  \texttt{Algebra.IsStandardSmoothOfRelativeDimension}\,\(n\)\,\(R\)\,\(S\),
  the algebra-side packaging of Stacks tag \texttt{00T7}(2)) and that
  the localisation \(S_M\) is non-trivial as a ring. Then the rank of
  \(\Omega_{S_M / R}\) over \(S_M\) equals \(n\).
\end{lemma}

\begin{proof}
  Mathlib's
  \texttt{Algebra.IsStandardSmoothOfRelativeDimension.\linebreak rank\_kaehlerDifferential}
  asserts \(\mathrm{rank}_S(\Omega_{S/R}) = n\) under the
  standard-smooth hypothesis. The localisation map
  \(\Omega_{S/R} \to \Omega_{S_M/R}\) is an
  \texttt{IsLocalizedModule} along \(S \to S_M\) (Stacks tag
  \texttt{02JK}; cf.\ the proof of
  \cref{lem:module_free_kaehler_localization}), and Mathlib's
  \texttt{Module.lift\_rank\_of\_isLocalizedModule\_of\_free} transports
  module rank through an \texttt{IsLocalizedModule} up to
  \(\mathrm{Cardinal.lift}\). Since \(n\) is a natural number both
  ranks equal \(n\), so
  \(\mathrm{rank}_{S_M}(\Omega_{S_M / R}) = n\) as required.
\end{proof}

\subsection{Stage 6 sub-decomposition: 00OE + 02JK chain (iter-198)}
\label{subsec:stage6_subgap_decomposition}

Stage 6 is the residual gap inside
\cref{lem:smooth_to_regular_local_ring}'s six-stage chain: given Stages
1--5b (a standard-smooth presentation of an affine chart
\(R \to S\) at a point \(z\), with \(\Omega_{S/R}\) free of rank \(n\)
and \(\Omega_{S_\mathfrak q / R}\) free of rank \(n\) on the stalk
\(S_\mathfrak q = \mathcal O_{X, z}\)), one must conclude that
\(S_\mathfrak q\) is a regular local ring. The Stacks proof of tag
\texttt{00TT} ultimately reduces to two ingredients which Mathlib does
not yet ship at commit \texttt{b80f227}: (6.A) the smooth-algebra
Krull-dimension formula (Stacks tag \texttt{00OE}, in the algebra
chapter as Lemma~\texttt{lemma-dimension-at-a-point-finite-type-field}),
and (6.B) the cotangent / K\"ahler-differential bridge over a field
base (Stacks tag \texttt{02JK}, in the algebra chapter as
Lemma~\texttt{lemma-differential-seq} plus the separable-residue-field
refinement of Lemma~\texttt{lemma-separable-smooth}). The two
ingredients combine (6.C) into the regularity conclusion via
\(\dim S_\mathfrak q = \mathrm{rank}_{S_\mathfrak q}\Omega_{S_\mathfrak
q / R} = \dim_{\kappa(\mathfrak q)} \mathfrak m / \mathfrak m^2\),
which is the cotangent-space characterisation of regularity that
Mathlib's \texttt{IsRegularLocalRing.iff\_finrank\_cotangentSpace}
already exposes. The next three lemmas package the three sub-gaps as
named iter-199+ prover targets.

\begin{lemma}
  [Stage 6.A --- Smooth-algebra Krull-dimension formula (Stacks 00OE)]
  \label{lem:smooth_algebra_krull_dim_formula}
  \lean{Algebra.IsStandardSmoothOfRelativeDimension.ringKrullDim_localization_eq_relativeDimension}
  \uses{lem:rank_kaehler_localization_eq_relative_dim,
        lem:ringKrullDim_localization_atMaximal_mvPolynomial,
        lem:ringKrullDim_quotient_localization_mvPolynomial_regular,
        lem:matsumura_isRegular_of_linearIndependent_cotangent,
        lem:exists_submersivePresentation_of_standardSmoothReldim,
        lem:submersive_relation_cotangent_linearIndependent_localized}
  % SOURCE: [Stacks Project], lemma-dimension-at-a-point-finite-type-field
  % (canonically known as tag 00OE; tag-to-label cross-check not in the
  % local `tags/tags` file, see iter-198 writer report note) (read from
  % references/stacks-algebra.tex, L28206--28217 in section
  % "Dimension theory of finitely generated $k$-algebras").
  % SOURCE QUOTE: "Let $k$ be a field. Let $S$ be a finite type $k$
  % algebra. Let $X = \Spec(S)$. Let $\mathfrak p \subset S$ be a prime
  % ideal, and let $x \in X$ be the corresponding point. Then we have
  % $\dim_x(X) = \dim(S_{\mathfrak p}) + \text{trdeg}_k\
  % \kappa(\mathfrak p)$."
  \textit{Source: [Stacks Project], tag~00OE (\(=\) \texttt{lemma-dimension-at-a-point-finite-type-field}).}
  Let \(k\) be a field, let \(S\) be a finite-type \(k\)-algebra, and let
  \(\mathfrak q \subset S\) be a prime ideal corresponding to a point
  \(x \in \mathrm{Spec}(S)\). Then
  \[
    \dim_x(\mathrm{Spec}\,S)
    \;=\; \dim S_{\mathfrak q} + \mathrm{trdeg}_k\,\kappa(\mathfrak q).
  \]
  In particular, if \(\mathfrak q\) is a maximal ideal (so \(x\) is a
  closed point) and \(k\) is algebraically closed, then
  \(\kappa(\mathfrak q) = k\) (Nullstellensatz), so
  \(\mathrm{trdeg}_k\,\kappa(\mathfrak q) = 0\) and
  \(\dim S_{\mathfrak q} = \dim_x(\mathrm{Spec}\,S)\).

  \paragraph{Geometric application.} Applied to the iter-193 Stage~3
  setup --- \(k = \bar k\), \(S = \Gamma(X.\mathrm{left}, V)\) a
  standard-smooth \(\bar k\)-algebra of relative dimension \(n\) (Stage
  3), and \(\mathfrak q\) the prime defining the point \(z \in V\) --- the
  formula yields \(\dim(\mathcal O_{X, z}) = \dim S_{\mathfrak q} = n -
  \mathrm{trdeg}_{\bar k}\,\kappa(z)\). For \(z\) a closed point this is
  simply \(n\); for the generic point of a codim-\(1\) subvariety
  (\(\mathrm{trdeg} = n - 1\)) it is \(1\).
\end{lemma}

% SOURCE QUOTE PROOF: "By Lemma \ref{lemma-dimension-prime-polynomial-ring}
% we know that $r = \text{trdeg}_k\ \kappa(\mathfrak p)$ is equal to the
% dimension of $V(\mathfrak p)$. Pick any maximal chain of primes
% $\mathfrak p \subset \mathfrak p_1 \subset \ldots \subset \mathfrak p_r$
% starting with $\mathfrak p$ in $S$. This has length $r$ by Lemma
% \ref{lemma-dimension-spell-it-out}. Let $\mathfrak q_j$, $j \in J$ be
% the minimal primes of $S$ which are contained in $\mathfrak p$. These
% correspond $1-1$ to minimal primes in $S_{\mathfrak p}$ via the rule
% $\mathfrak q_j \mapsto \mathfrak q_jS_{\mathfrak p}$. By Lemma
% \ref{lemma-dimension-at-a-point-finite-type-over-field} we know that
% $\dim_x(X)$ is equal to the maximum of the dimensions of the rings
% $S/\mathfrak q_j$. For each $j$ pick a maximal chain of primes
% $\mathfrak q_j \subset \mathfrak p'_1 \subset \ldots \subset
% \mathfrak p'_{s(j)} = \mathfrak p$. Then $\dim(S_{\mathfrak p}) =
% \max_{j \in J} s(j)$. Now, each chain
% $\mathfrak q_j \subset \mathfrak p'_1 \subset \ldots \subset
% \mathfrak p'_{s(j)} = \mathfrak p \subset \mathfrak p_1 \subset \ldots
% \subset \mathfrak p_r$ is a maximal chain in $S/\mathfrak q_j$, and by
% what was said before we have $\dim_x(X) = \max_{j \in J} r + s(j)$.
% The lemma follows."
\begin{proof}
  \uses{lem:rank_kaehler_localization_eq_relative_dim,
        lem:ringKrullDim_localization_atMaximal_mvPolynomial,
        lem:ringKrullDim_quotient_localization_mvPolynomial_regular,
        lem:matsumura_isRegular_of_linearIndependent_cotangent,
        lem:exists_submersivePresentation_of_standardSmoothReldim,
        lem:submersive_relation_cotangent_linearIndependent_localized}
  This is exactly Stacks tag~\texttt{00OE}
  (\texttt{lemma-dimension-at-a-point-finite-type-field}); the proof is a
  citation of the dimension theory of finitely generated algebras over a
  field. By Noether normalisation the finite-type \(k\)-algebra \(S\) is
  finite over a polynomial subalgebra, so the going-up / catenary
  dimension theory of such algebras applies. Write
  \(r = \mathrm{trdeg}_k\,\kappa(\mathfrak q)\); this transcendence
  degree is exactly the dimension of the closed subvariety
  \(\overline{\{x\}} = V(\mathfrak q)\). A maximal chain of primes
  through \(\mathfrak q\) therefore splits into a
  length-\(\dim S_{\mathfrak q}\) part below \(\mathfrak q\) (recording
  the local codimension, i.e.\ the height of \(\mathfrak q\) in the
  relevant irreducible component) and a length-\(r\) part above it
  (recording the dimension of \(\overline{\{x\}}\)). Taking the maximum
  over the minimal primes contained in \(\mathfrak q\) yields
  \[
    \dim_x(\mathrm{Spec}\,S) = \dim S_{\mathfrak q} + r
    = \dim S_{\mathfrak q} + \mathrm{trdeg}_k\,\kappa(\mathfrak q),
  \]
  which is the asserted formula.

  For the geometric specialisation matching the Lean target, take
  \(k = \bar k\) algebraically closed and \(S\) standard smooth of
  relative dimension \(n\), so that \(\dim_x(\mathrm{Spec}\,S) = n\) via
  \cref{lem:rank_kaehler_localization_eq_relative_dim} (the free
  K\"ahler module \(\Omega_{S_{\mathfrak q}/R}\) has rank \(n\), which
  pins the relative dimension). The formula then gives
  \(\dim S_{\mathfrak q} = n - \mathrm{trdeg}_{\bar k}\,\kappa(\mathfrak
  q)\). At a closed point \(\mathfrak q\) is maximal, so by the
  Nullstellensatz \(\kappa(\mathfrak q) = \bar k\) and
  \(\mathrm{trdeg}_{\bar k}\,\kappa(\mathfrak q) = 0\); hence
  \(\dim S_{\mathfrak q} = \dim_x(\mathrm{Spec}\,S) = n\).
\end{proof}

\begin{lemma}
  [Stage 6.B --- Cotangent \(\leftrightarrow\) K\"ahler over a field (Stacks 02JK)]
  \label{lem:cotangent_kahler_over_field}
  \lean{AlgebraicGeometry.Scheme.cotangent_iso_residue_tensor_kaehler_of_formallySmooth_residue}
  % \lean{...} updated iter-199 review: iter-199 prover landed the
  % closed-point iso as a project-local axiom-clean helper (private
  % theorem in AlgebraicGeometry.Scheme namespace). Companion siblings:
  %   cotangent_iso_maximalIdeal_residue_tensor_kaehler_of_formallySmooth_residue
  %   finrank_cotangentSpace_of_formallySmooth_residue
  %   finrank_cotangentSpace_of_bijective_algebraMap_residue (closed-point bundled)
  \uses{lem:smooth_algebra_krull_dim_formula}
  % SOURCE: [Stacks Project], lemma-differential-seq (canonically
  % known as the right-exact conormal portion of tag 02JK; the
  % short-exact / split refinement is the proof step inside
  % lemma-separable-smooth) (read from
  % references/stacks-algebra.tex, L34087--34104 and L38724--38766).
  % SOURCE QUOTE: "In diagram (\ref{equation-functorial-omega}),
  % suppose that $S \to S'$ is surjective with kernel $I \subset S$, and
  % assume that $R' = R$. Then there is a canonical exact sequence of
  % $S'$-modules
  % $$I/I^2 \longrightarrow \Omega_{S/R} \otimes_S S' \longrightarrow
  % \Omega_{S'/R} \longrightarrow 0$$
  % The leftmost map is characterized by the rule that $f \in I$ maps
  % to $\text{d}f \otimes 1$."
  % SOURCE QUOTE: "Hence we get
  % $$\dim_{\kappa} \Omega_{R/k} \otimes_R \kappa = \dim_\kappa
  % \mathfrak m/\mathfrak m^2 + \text{trdeg}_k (\kappa) \geq \dim R +
  % \text{trdeg}_k (\kappa) = \dim_{\mathfrak q} S$$
  % (see Lemma \ref{lemma-dimension-at-a-point-finite-type-field} for
  % the last equality) with equality if and only if $R$ is regular."
  \textit{Source: [Stacks Project], tag~02JK (right-exact half = \texttt{lemma-differential-seq}; separable-residue refinement inside \texttt{lemma-separable-smooth}).}
  Let \(k\) be a field and let \(R\) be a Noetherian local
  \(k\)-algebra with maximal ideal \(\mathfrak m\) and residue field
  \(\kappa = R/\mathfrak m\). The canonical exact sequence
  \[
    \mathfrak m / \mathfrak m^2
    \;\longrightarrow\; \Omega_{R/k} \otimes_R \kappa
    \;\longrightarrow\; \Omega_{\kappa / k}
    \;\longrightarrow\; 0
  \]
  identifies the cotangent space \(\mathfrak m / \mathfrak m^2\) (a
  \(\kappa\)-vector space) with the kernel of the right map. When
  \(\kappa / k\) is separable (in particular when \(k\) is
  algebraically closed and \(\kappa = k\), so \(\Omega_{\kappa / k}
  = 0\)) the leftmost map is injective and the sequence is short exact
  --- in fact split short exact when there is a \(k\)-algebra section
  \(\kappa \to R\) (e.g.\ at a closed point of a \(\bar k\)-variety).
  In particular, over an algebraically closed field \(k = \bar k\) and
  at a closed point \(\mathfrak m\),
  \[
    \Omega_{R/k} \otimes_R \kappa \;\simeq\; \mathfrak m / \mathfrak m^2
  \]
  as \(\kappa\)-vector spaces; hence
  \(\dim_\kappa(\Omega_{R/k} \otimes_R \kappa)
   = \dim_\kappa \mathfrak m / \mathfrak m^2\).
\end{lemma}

The Stage 6.B cotangent--K\"ahler chain is realised on the Lean side by the
following substrate lemmas, which compute the residue-field tensor finrank and
promote the cotangent isomorphism of \cref{lem:cotangent_kahler_over_field} into
the cotangent-space finrank consumed by the regularity criterion.

\begin{lemma}\leanok
[Stage 6.B (RHS) --- residue-tensor finrank from free rank]
  \label{lem:finrank_residue_kaehler_of_free_rank}
  \lean{AlgebraicGeometry.Scheme.finrank_residueField_tensor_kaehlerDifferential_of_free_rank_eq}
  Let \(R \to S_{\mathfrak q}\) be an algebra with \(S_{\mathfrak q}\) a nontrivial
  local ring whose K\"ahler differentials \(\Omega_{S_{\mathfrak q}/R}\) are
  \(S_{\mathfrak q}\)-free of rank \(n\). Then the residue-field base change
  \(\kappa \otimes_{S_{\mathfrak q}} \Omega_{S_{\mathfrak q}/R}\) has
  \(\kappa\)-dimension \(n\), where \(\kappa\) is the residue field of \(S_{\mathfrak q}\).
\end{lemma}

\begin{proof}
  Proved directly in Lean: base change preserves rank, and finite rank equals the
  natural number \(n\). A sub-step of
  \cref{lem:finrank_cotangentSpace_of_formallySmooth_residue}.
\end{proof}

\begin{lemma}[Stage 6.B (RHS) --- residue-tensor finrank for a standard-smooth algebra]
  \label{lem:finrank_residue_kaehler_of_standardSmooth_reldim}
  \lean{AlgebraicGeometry.Scheme.finrank_residueField_tensor_kaehlerDifferential_of_isStandardSmoothOfRelativeDimension}
  \uses{lem:finrank_residue_kaehler_of_free_rank}
  If \(S_{\mathfrak q}\) is a nontrivial local \(R\)-algebra that is standard smooth
  of relative dimension \(n\), then the residue-field base change of its K\"ahler
  differentials has \(\kappa\)-dimension \(n\). This is the standard-smooth-tied
  packaging of \cref{lem:finrank_residue_kaehler_of_free_rank}, discharging the
  freeness and rank hypotheses directly from the relative-dimension instance.
\end{lemma}

\begin{proof}
  Proved directly in Lean: the standard-smooth instance supplies freeness and
  \(\mathrm{rank}\,\Omega_{S_{\mathfrak q}/R} = n\); apply
  \cref{lem:finrank_residue_kaehler_of_free_rank}.
\end{proof}

\begin{lemma}[Stage 6.B --- cotangent isomorphism, maximal-ideal form]
  \label{lem:cotangent_iso_maximalIdeal_residue_kaehler}
  \lean{AlgebraicGeometry.Scheme.cotangent_iso_maximalIdeal_residue_tensor_kaehler_of_formallySmooth_residue}
  \uses{lem:cotangent_kahler_over_field}
  Under the closed-point hypotheses of \cref{lem:cotangent_kahler_over_field}
  (\(S_{\mathfrak q}\) formally smooth over \(R\), residue field \(\kappa\) formally
  smooth over \(R\), and \(\Omega_{\kappa/R} = 0\)), there is an
  \(S_{\mathfrak q}\)-linear isomorphism between the cotangent module of the maximal
  ideal \(\mathfrak m\) and the residue-field base change
  \(\kappa \otimes_{S_{\mathfrak q}} \Omega_{S_{\mathfrak q}/R}\). This restates the
  kernel-side cotangent isomorphism of \cref{lem:cotangent_kahler_over_field} with the
  canonical maximal-ideal cotangent space as its domain.
\end{lemma}

\begin{proof}
  Proved directly in Lean: transports the isomorphism of
  \cref{lem:cotangent_kahler_over_field} along the equality
  \(\ker(S_{\mathfrak q} \to \kappa) = \mathfrak m\).
\end{proof}

\begin{lemma}\leanok
[Stage 6.B --- cotangent-space finrank, formally-smooth residue form]
  \label{lem:finrank_cotangentSpace_of_formallySmooth_residue}
  \lean{AlgebraicGeometry.Scheme.finrank_cotangentSpace_of_formallySmooth_residue}
  \uses{lem:cotangent_iso_maximalIdeal_residue_kaehler, lem:finrank_residue_kaehler_of_free_rank}
  With the closed-point hypotheses of
  \cref{lem:cotangent_iso_maximalIdeal_residue_kaehler} and
  \(\Omega_{S_{\mathfrak q}/R}\) free of \(S_{\mathfrak q}\)-rank \(n\), the cotangent
  space \(\mathfrak m / \mathfrak m^2\) has \(\kappa\)-dimension \(n\).
\end{lemma}

\begin{proof}
  Proved directly in Lean: combine the \(\kappa\)-promoted cotangent isomorphism
  \cref{lem:cotangent_iso_maximalIdeal_residue_kaehler} with the residue-tensor
  finrank \cref{lem:finrank_residue_kaehler_of_free_rank}.
\end{proof}

\begin{lemma}\leanok
[Stage 6.B --- cotangent-space finrank at a $\bar k$-rational point]
  \label{lem:finrank_cotangentSpace_of_bijective_residue}
  \lean{AlgebraicGeometry.Scheme.finrank_cotangentSpace_of_bijective_algebraMap_residue}
  \uses{lem:finrank_cotangentSpace_of_formallySmooth_residue}
  Bundled closed-point form of
  \cref{lem:finrank_cotangentSpace_of_formallySmooth_residue}: if \(S_{\mathfrak q}\)
  is formally smooth over \(R\), \(\Omega_{S_{\mathfrak q}/R}\) is free of rank \(n\),
  and the structure map \(R \to \kappa\) to the residue field is bijective (as at a
  \(\bar k\)-rational closed point, by the Nullstellensatz), then
  \(\dim_\kappa \mathfrak m/\mathfrak m^2 = n\).
\end{lemma}

\begin{proof}
  Proved directly in Lean: bijectivity upgrades the residue field to formally smooth
  with vanishing differentials, reducing to
  \cref{lem:finrank_cotangentSpace_of_formallySmooth_residue}.
\end{proof}

\begin{lemma}
  [Stage 6.C --- Assembly: smooth standard-smooth stalk is regular local]
  \label{lem:stage6_regular_stalk_assembly}
  \lean{AlgebraicGeometry.TODO.stage6_regular_stalk_assembly}
  % NOTE (iter-270 dag-walker): \lean{...} pinned to the TODO placeholder
  % decl name. This 6.C assembly has NO separately-pinnable Lean
  % declaration: its content is the in-body closure pattern of
  % \texttt{isRegularLocalRing\_stalk\_of\_smooth} (pinned by
  % \cref{lem:smooth_to_regular_local_ring}), so it cannot re-pin that same
  % decl. The TODO pin gives leandag a resolvable \lean{...} target so the
  % node carries a name; it is NOT a claim that a kernel-clean decl exists.
  % NOTE (iter-199): the Lean target for this 6.C assembly is the body
  % of \texttt{AlgebraicGeometry.Scheme.isRegularLocalRing\_stalk\_of\_smooth}
  % itself (see \cref{lem:smooth_to_regular_local_ring}); the iter-198
  % blueprint named a planned \texttt{...\_aux} extraction that has not
  % been created in the Lean file. Treat the lemma below as the in-body
  % closure pattern that the existing helper consumes once Stages 6.A and
  % 6.B land, rather than as a separately-pinned Lean declaration.
  \uses{lem:rank_kaehler_localization_eq_relative_dim,
        lem:smooth_algebra_krull_dim_formula,
        lem:cotangent_kahler_over_field}
  Let \(\bar k\) be an algebraically closed field, let
  \(R \to S\) be a standard-smooth \(\bar k\)-algebra presentation of
  relative dimension \(n\) on an affine chart \(\Spec S \subseteq X\),
  and let \(z \in \Spec S\) be a point corresponding to a prime
  \(\mathfrak q \subset S\). Then the stalk \(S_{\mathfrak q}\) is a
  regular local ring of Krull dimension
  \(n - \mathrm{trdeg}_{\bar k}\,\kappa(\mathfrak q)\). In the
  closed-point case (\(\kappa(\mathfrak q) = \bar k\),
  \(\mathrm{trdeg} = 0\)) the dimension equals \(n\); in the
  codim-\(1\) generic-point case (\(\mathrm{trdeg} = n - 1\)) the
  dimension equals \(1\).
\end{lemma}

\begin{proof}
  Combine the three sub-lemmas:
  \begin{enumerate}
    \item By \cref{lem:rank_kaehler_localization_eq_relative_dim} (Stage~5b)
      the localised K\"ahler module \(\Omega_{S_{\mathfrak q} / R}\)
      is free of rank \(n\) over \(S_{\mathfrak q}\).
    \item By \cref{lem:smooth_algebra_krull_dim_formula} (Stage 6.A,
      Stacks 00OE) the Krull dimension of \(S_{\mathfrak q}\) equals
      \(\dim_z \Spec S - \mathrm{trdeg}_{\bar k}\,\kappa(\mathfrak q)\),
      which for a standard-smooth presentation of relative dimension
      \(n\) equals \(n - \mathrm{trdeg}_{\bar k}\,\kappa(\mathfrak q)\).
      (For the standard-smooth presentation the dimension
      \(\dim_z \Spec S = n\) follows from the global-complete-intersection
      structure: \(n\) variables modulo \(n - \dim\) equations, see
      Stacks \texttt{lemma-relative-global-complete-intersection-smooth}.)
    \item By \cref{lem:cotangent_kahler_over_field} (Stage 6.B,
      Stacks 02JK) the residue-field tensor \(\Omega_{S_{\mathfrak q}/R}
      \otimes_{S_{\mathfrak q}} \kappa(\mathfrak q)\) is isomorphic to
      \(\mathfrak m / \mathfrak m^2\) as \(\kappa(\mathfrak q)\)-vector
      spaces. Combined with (1), \(\dim_{\kappa(\mathfrak q)} \mathfrak
      m / \mathfrak m^2 = n\) at the closed-point case (and in general
      \(n - \mathrm{trdeg}_{\bar k}\,\kappa(\mathfrak q)\) by base-change
      of the standard-smooth presentation to the closed-point fibre).
    \item Compare (2) and (3): both \(\dim S_{\mathfrak q}\) and
      \(\dim_{\kappa(\mathfrak q)} \mathfrak m / \mathfrak m^2\) equal
      \(n - \mathrm{trdeg}_{\bar k}\,\kappa(\mathfrak q)\); the
      regularity criterion for Noetherian local rings (the equality
      \(\dim_\kappa \mathfrak m / \mathfrak m^2 = \dim R\) is the
      defining condition for a regular local ring, Stacks
      \texttt{definition-regular-local} \(=\) Mathlib
      \texttt{IsRegularLocalRing}) then gives \(S_{\mathfrak q}\)
      regular. The scheme-level transfer
      \(S_{\mathfrak q} = \mathcal O_{X, z}\) (Stage 3
      \texttt{IsLocalization.AtPrime} on the affine chart) completes
      the conclusion.
  \end{enumerate}
\end{proof}

\paragraph{Mathlib API state (commit \texttt{b80f227}).} The following
pieces are needed to formalise Stages 6.A--6.C; status as audited iter-198.

\begin{itemize}
  \item \texttt{Algebra.IsStandardSmoothOfRelativeDimension} --- EXISTS
    in Mathlib (\texttt{Mathlib/RingTheory/Smooth/StandardSmooth.lean});
    relative-dimension witness already consumed by Stage 5b
    (\cref{lem:rank_kaehler_localization_eq_relative_dim}).
  \item \texttt{Algebra.FormallySmooth} / \texttt{Algebra.Smooth} ---
    EXISTS (\texttt{Mathlib/RingTheory/Smooth/Basic.lean}); used in
    Stages 1--2 to deduce standard-smooth presentations.
  \item \texttt{ringKrullDim} / \texttt{Ring.KrullDim} --- EXISTS
    (\texttt{Mathlib/RingTheory/KrullDimension/Basic.lean}); supplies
    the Krull-dimension side. The link
    \texttt{ringKrullDim S\_q = m.height} is supplied by
    \texttt{IsLocalization.AtPrime.ringKrullDim\_eq\_height}
    (\texttt{Mathlib/RingTheory/Ideal/Height.lean}). EXISTS.
  \item \texttt{Algebra.IsStandardSmoothOfRelativeDimension.\linebreak ringKrullDim\_localization\_eq\_relativeDimension}
    (the project's name for Stage 6.A, Stacks 00OE) --- MISSING in
    Mathlib at \texttt{b80f227} as a packaged form. Iter-200 closed the
    polynomial-ring side axiom-clean via 7 \texttt{private} substrate
    declarations in \texttt{CodimOneExtension.lean} L688--L800 (see
    \ref{subsec:stage6_iib_substrate_iter200}): the height-of-maximal-ideal
    formula \texttt{MvPolynomial.maximalIdeal\_height\_eq\_card / \_eq\_natCard}
    on \texttt{MvPolynomial} \(k\) is in hand by induction on \texttt{Fin n}
    via \texttt{Polynomial.height\_eq\_height\_add\_one}; the capstone
    \texttt{ringKrullDim\_localization\_atMaximal\_MvPolynomial} composes
    Steps 1+2 into a single consumer call. Residual at iter-201 is
    Step 3: a Jacobian-criterion regular-sequence witness for the
    relations of the smooth-algebra presentation (Stacks 00SW / 00OW,
    \(\sim 30-60\) LOC project-side; recipe in
    \ref{subsec:stage6_iib_substrate_iter200} below).
  \item \texttt{KaehlerDifferential.exact\_mapBaseChange} (the
    conormal-sequence right-exactness, \(=\) Stacks
    \texttt{lemma-differential-seq}) --- EXISTS in Mathlib at
    \texttt{Mathlib/RingTheory/Kaehler/CotangentComplex.lean} (some
    naming variant of \texttt{conormal\_seq\_exact} /
    \texttt{KaehlerDifferential.map\_surjective}). Verify the exact name
    at \texttt{b80f227} before pinning.
  \item \texttt{KaehlerDifferential.cotangent\_iso\_residue\_tensor\_kaehler}
    (the project's name for Stage 6.B, Stacks 02JK at a separable
    residue point) --- MISSING in Mathlib at \texttt{b80f227} in this
    packaged form; pieces are present (conormal sequence + \(\Omega_{\kappa
    / k} = 0\) for \(\kappa = k\) algebraically closed), but the
    \texttt{IsLocalRing.CotangentSpace} \(\leftrightarrow\)
    \texttt{KaehlerDifferential} bridge is not assembled. NEEDS-BRIDGE;
    project build of \(\sim\)100--200~LOC.
  \item \texttt{IsRegularLocalRing.iff\_finrank\_cotangentSpace} ---
    EXISTS (\texttt{Mathlib/RingTheory/RegularLocal/Basic.lean}); this
    is the cotangent-space characterisation of regularity that
    Stage 6.C feeds into. Consumes \texttt{Module.finrank} of
    \texttt{IsLocalRing.CotangentSpace} \(=\) \(\dim_\kappa \mathfrak m
    / \mathfrak m^2\).
  \item \texttt{Algebra.IsSmooth.krullDim\_stalk} or any direct
    smooth-algebra \(\Rightarrow\) regular-stalk shortcut --- MISSING
    at \texttt{b80f227}. The standard scheme-level packaging of Stacks
    00TT is not in Mathlib; the iter-191--193 staged decomposition is
    the project-side substitute.
  \item \texttt{IsRegularLocalRing.localization\_isRegularLocalRing}
    (preservation of regularity under localisation at a prime;
    Stacks \texttt{tag 00OF}) --- \textbf{MISSING at \texttt{b80f227}}
    (corrected iter-202: iter-201 Lane COE prover scout confirmed via
    \texttt{grep} of \texttt{Mathlib/RingTheory/RegularLocalRing/}
    that only \texttt{Defs.lean} is present). The
    closed-point-to-arbitrary-point alternative route through this
    lemma is therefore also gated. The main Step A1 route does not
    need this lemma (see
    \cref{subsec:stage6_iib_substrate_iter200}).
\end{itemize}

\paragraph{Cascade to consumers.} Closure of Stage 6 (i.e.\ landing
\cref{lem:stage6_regular_stalk_assembly} axiom-clean) collapses the
following residual sorries in
\texttt{AlgebraicJacobian/Albanese/CodimOneExtension.lean}:
\begin{itemize}
  \item \texttt{isRegularLocalRing\_stalk\_of\_smooth} (L434--L526) ---
    its final \texttt{sorry} on L526 is filled by chaining the Stage 5a/5b
    helpers with the new Stage 6.A/6.B/6.C bridges, exactly as the
    proof of \cref{lem:stage6_regular_stalk_assembly} sketches. Closure
    is automatic once the three sub-lemmas land axiom-clean.
  \item \texttt{smooth\_codim\_one\_maximalIdeal\_isPrincipal\_and\_ne\_bot}
    (L547--L599) --- already axiom-clean modulo the named Stacks-00TT
    gap; the regularity half is delegated to
    \texttt{isRegularLocalRing\_stalk\_of\_smooth}, so closure
    propagates upward without further work.
  \item \texttt{localRing\_dvr\_of\_codim\_one} (L615--L644) ---
    likewise upgrades from gated to axiom-clean automatically: it
    consumes \texttt{smooth\_codim\_one\_maximalIdeal\_isPrincipal\_and\_ne\_bot}
    and then closes via \texttt{IsDiscreteValuationRing.TFAE}.
  \item \texttt{extend\_of\_codimOneFree\_of\_smooth} (the body
    around L703--L723) --- Step~1 of
    \cref{thm:codim_one_extension}'s proof consumes
    \texttt{localRing\_dvr\_of\_codim\_one} via the valuative
    criterion at codim-\(1\) points, so Stage 6 closure unblocks
    Step~1 entirely. Step~2 (the depth-\(\geq 2\) local-cohomology
    vanishing, Stacks tag \texttt{0AVF}) remains an independent
    Mathlib gap and is \emph{not} unblocked by Stage 6.
  \item \texttt{indeterminacy\_pure\_codim\_one\_into\_grpScheme}
    (L752--L798, the Lean body of
    \cref{lem:milne_codim1_indeterminacy}) --- not directly gated on
    Stage 6 (its substantive Mathlib gap is the codim-\(1\)
    pole-divisor / diagonal intersection lemma, Sub-step~4 of
    Milne's proof). Stage 6 closure does \emph{not} unblock this
    lemma; it is recorded here only for completeness of the
    chapter's dependency picture.
\end{itemize}

The two remaining post-Stage-6 gaps in the chapter are therefore:
(a) the Step~2 depth-\(\geq 2\) local-cohomology vanishing
(Stacks~\texttt{0AVF}) consumed only by
\cref{thm:codim_one_extension}, and (b) the
function-field-pullback machinery for \texttt{Scheme.RationalMap}
consumed by \cref{thm:weil_divisor_obstruction} and Sub-step~3 of
\cref{lem:milne_codim1_indeterminacy}. Neither overlaps with the
Stage 6 work, so the three sub-lemmas above are a self-contained
mathlib-analogist target.

\subsection{Stage 6.A iter-200 substrate: Stacks 00OE 3-step decomposition}
\label{subsec:stage6_iib_substrate_iter200}
\textit{Source: \texttt{coe-stacks00oe} mathlib-analogist
cross-domain-inspiration recipe (iter-200, persistent
\texttt{analogies/coe-stacks00oe.md}); 7 axiom-clean
\texttt{private} substrate theorems landed in
\texttt{CodimOneExtension.lean} L688--L800 under
\texttt{namespace AlgebraicGeometry.Scheme}.}

The iter-200 Lane COE prover phase decomposed Stage 6.A
(\ref{lem:smooth_algebra_krull_dim_formula}, Stacks 00OE) into a
3-step Mathlib-chain. Steps 1 and 2 landed axiom-clean as
project-local \texttt{private} theorems; Step 3 landed in an
additive parameterised form awaiting the Jacobian-regular-sequence
witness for the relations of the smooth-algebra presentation.

\paragraph{Step 1 --- Krull-dim \(\Rightarrow\) height (trivial,
\(\sim 1\) line).} Re-export of Mathlib's
\(\texttt{IsLocalization.AtPrime.ringKrullDim\_eq\_height}\) under
the project-local name
\(\texttt{ringKrullDim\_localization\_eq\_height\_atPrime}\); reduces
the localised Krull-dim goal to a height-of-maximal-ideal goal.

\paragraph{Step 2 --- Polynomial-ring maximal-ideal height
(\(\sim 60\) LOC).} For \(\texttt{MvPolynomial}\,(\texttt{Fin}\,n)\,k\)
with \(k\) a field, every maximal ideal has height \emph{exactly}
\(n\). Lower bound:
\(\texttt{MvPolynomial.maximalIdeal\_height\_ge\_card\_of\_field}\),
induction on \texttt{Fin n} via \(\texttt{MvPolynomial.finSuccEquiv}\)
+ Mathlib's \(\texttt{Polynomial.height\_eq\_height\_add\_one}\) plus
the Jacobson contraction
\(\texttt{Polynomial.isMaximal\_comap\_C\_of\_isJacobsonRing}\). Upper
bound:
\(\texttt{MvPolynomial.maximalIdeal\_height\_le\_natCard\_of\_field}\)
via \(\texttt{Ideal.height\_le\_ringKrullDim\_of\_ne\_top}\) +
\(\texttt{MvPolynomial.ringKrullDim\_of\_isNoetherianRing}\) +
\(\texttt{ringKrullDim\_eq\_zero\_of\_field}\). Combined: the
\texttt{Fin n}-form
\(\texttt{MvPolynomial.maximalIdeal\_height\_eq\_card}\) and the
general \texttt{[Finite ι]}-form
\(\texttt{MvPolynomial.maximalIdeal\_height\_eq\_natCard}\) via
\(\texttt{MvPolynomial.renameEquiv}\) +
\(\texttt{RingEquiv.height\_map}\). Capstone packaging the
Steps 1+2 composition for any localisation at a maximal ideal:
\(\texttt{ringKrullDim\_localization\_atMaximal\_MvPolynomial}\).

\paragraph{Step 3 --- Drop by regular-sequence length (additive form
in hand, witness pending).} Step 3 generally:
\(\texttt{ringKrullDim}\,(R'_{m'} / (f_j)_{j \in \sigma}) +
 |\sigma| = \texttt{ringKrullDim}\,R'_{m'}\). The additive form is
landed as
\(\texttt{ringKrullDim\_quotient\_add\_eq\_of\_regular\_sequence}\)
via Mathlib's
\(\texttt{ringKrullDim\_add\_length\_eq\_ringKrullDim\_of\_isRegular}\),
parameterised over a hypothesis
\(\texttt{RingTheory.Sequence.IsRegular}\,R'_{m'}\,\texttt{rs}\)
that is the iter-201+ residual substrate obligation. Note:
\(\texttt{WithBot}\,\mathbb N_\infty\) carries no \texttt{HSub}
instance, so the additive form is the natural API; a subtraction
form is not available.

\paragraph{Iter-201+ Jacobian-regular-sequence witness recipe.} The
residual substrate of Stage 6.A is the Jacobian-criterion
regular-sequence witness for the relations of an
\(\texttt{Algebra.SubmersivePresentation}\).

\textbf{API choice} (per iter-201 \texttt{progress-critic route201}
flag): the project-local target predicate is Mathlib's
\(\texttt{RingTheory.Sequence.IsRegular}\), \emph{not}
\(\texttt{IsWeaklyRegular}\). Reason: the iter-200 substrate
\(\texttt{ringKrullDim\_quotient\_add\_eq\_of\_regular\_sequence}\)
(\texttt{CodimOneExtension.lean} L807) consumes the strong predicate
\(\texttt{IsRegular}\) (via Mathlib's
\(\texttt{ringKrullDim\_add\_length\_eq\_ringKrullDim\_of\_isRegular}\));
the localised stalk \(R'_{m'}\) is local, so by Mathlib's
\(\texttt{IsLocalRing.isRegular\_iff\_isWeaklyRegular\_of\_subset\_maximalIdeal}\)
the two predicates are equivalent under the standing hypothesis that
the relations lie in the maximal ideal --- which holds because
the Jacobian-criterion conclusion is precisely that the relations
generate (a subset of) the maximal ideal at the closed point. The
recipe builds \(\texttt{IsWeaklyRegular}\) first (it is the
naturally-arising predicate from the Koszul-complex argument) and
then promotes to \(\texttt{IsRegular}\) via the bridge lemma; the
final consumer call uses the \(\texttt{IsRegular}\) form.

The recipe (per the iter-200 prover handoff + the
\texttt{coe-stacks00oe} analogist + iter-201 \texttt{coe-stacks00sw}
analogist):
(i) extract the explicit
\(\texttt{Algebra.SubmersivePresentation}\,\bar k\,
 \Gamma(X.\texttt{left}, V)\) from
\(\texttt{Algebra.IsStandardSmoothOfRelativeDimension}\) (Stage 6 sub-gap
(i) substrate iter-198); (ii) cite Stacks 00SW / 00OW + Matsumura
Thm 14.2: the relations \((f_j)_{j \in \sigma}\) of a submersive
presentation form a Koszul-regular sequence at the closed point
because the Jacobian matrix is invertible modulo the maximal ideal
\(\Leftrightarrow\) the images of \(f_j\) in \(\mathfrak m_A /
\mathfrak m_A^2\) are \(\kappa\)-linearly independent
\(\Rightarrow\) (Matsumura 14.2 / Stacks 00NQ on a regular local
ring) the \(f_j\) form a regular sequence in \(A\); (iii) bundle
as a 3-step chain:
\begin{itemize}
  \item \textbf{A1 — Matsumura helper} (revised cost
    \(\sim 80\)--\(150\) LOC \emph{including cross-file imports}; the
    earlier \(\sim 30\)--\(50\) LOC estimate did not account for
    project-local Stacks 00NQ + RLR-preservation supply): on a regular
    local Noetherian ring \((A, \mathfrak m)\) of Krull dim \(n\), a
    sequence \(f_1, \dots, f_c \in \mathfrak m\) whose images in
    \(\mathfrak m / \mathfrak m^2\) are \(\kappa\)-linearly
    independent forms an \(\texttt{IsRegular}\) sequence in \(A\).
    Proof: induction on \(c\), dim-drops-by-one via Mathlib's
    \(\texttt{ringKrullDim\_quotient\_span\_singleton\_succ\_eq\_ringKrullDim}\)
    (in \texttt{Mathlib/RingTheory/KrullDimension/Regular.lean}),
    \(\texttt{IsSMulRegular}\) on the non-zero-divisor \(f_1\) via
    \(\texttt{IsRegularLocalRing} \Rightarrow \texttt{IsDomain}\) +
    \(f_1 \neq 0\) (since \(f_1 \notin \mathfrak m^2\)), then
    \(\texttt{isRegular\_cons\_iff}\) to descend.

    \textbf{Mathlib API state (iter-201 prover scout + iter-201
    \texttt{lean-auditor iter201} cross-file finding).} Two of the
    ingredients are \emph{absent} from Mathlib at \texttt{b80f227} but
    \emph{present project-locally} as private witnesses in
    \texttt{AlgebraicJacobian/Albanese/AuslanderBuchsbaum.lean}, all
    axiom-clean:
    \begin{itemize}
      \item \(\texttt{IsRegularLocalRing} \Rightarrow \texttt{IsDomain}\)
        (Stacks 00NQ) --- Mathlib has only the PID-converse instance
        \([\texttt{IsLocalRing R}]\,[\texttt{IsDomain R}]\,
         [\texttt{IsPrincipalIdealRing R}]\,
         \to \texttt{IsRegularLocalRing R}\). The forward direction is
        supplied project-side by
        \(\texttt{RingTheory.CohenMacaulay.isDomain\_of\_regularLocal}\)
        (\texttt{AuslanderBuchsbaum.lean} L2657, public as of
        iter-202; strong induction on
        \(\texttt{spanFinrank}\,\mathfrak m\) closed iter-186 modulo
        \(\texttt{notMem\_minimalPrimes\_of\_regularLocal\_succ}\),
        latter closed iter-201). The iter-202 Lane AB dispatch
        landed the \texttt{private}\(\to\)public promotion, so the
        name is reachable cross-file.
      \item \(A / (f_1) \to \texttt{IsRegularLocalRing}\)
        preservation under quotient by an element not in
        \(\mathfrak m^2\) --- absent in Mathlib (only the dim-drop
        \(\texttt{ringKrullDim\_quotient\_span\_singleton\_succ\_eq\_ringKrullDim}\)
        is present, no RLR-preservation companion). Supplied
        project-side by
        \(\texttt{RingTheory.CohenMacaulay.regularLocal\_quotient\_isRegularLocal\_of\_notMemSq}\)
        (\texttt{AuslanderBuchsbaum.lean} L2293, public as of
        iter-202). Same iter-202 Lane AB promotion path, landed.
    \end{itemize}
    Both project-local witnesses are now public and reachable
    cross-file (the iter-202 Lane AB dispatch removed
    \texttt{private}), so Step A1 closes in the projected
    \(\sim 30\)--\(50\) LOC via the recipe above; the cross-file
    blocker that fenced iter-202 is discharged.
    \textbf{Iter-203 Lane COE Step A1 prover recipe (full
    specification per iter-202 \texttt{progress-critic route202}
    must-fix).} For the iter-203 prover dispatch (iter-202 Lane AB
    landed the two promotions), the recipe is:

    \emph{Imports} (top of
    \texttt{AlgebraicJacobian/Albanese/CodimOneExtension.lean}, in
    addition to existing): \texttt{import
    AlgebraicJacobian.Albanese.AuslanderBuchsbaum} (already present).
    After the iter-202 Lane AB private removals, the qualified names
    \(\texttt{RingTheory.CohenMacaulay.isDomain\_of\_regularLocal}\) and
    \(\texttt{RingTheory.CohenMacaulay.regularLocal\_quotient\_isRegularLocal\_of\_notMemSq}\)
    become reachable cross-file with no new import line.

    \emph{Target declaration} (new \texttt{private theorem}, insert
    between iter-201 \texttt{ringKrullDim\_quotient\_localization\_MvPolynomial\_of\_regular}
    L924 and iter-198 \texttt{exists\_isStandardSmoothOfRelativeDimension\_of\_isStandardSmooth}
    L834+):
    \begin{verbatim}
    -- LANDED iter-203 (axiom-clean). The dimension is read off the
    -- cotangent-independence internally, so the explicit
    -- `(hn : ringKrullDim A = n)` parameter of the original recipe was
    -- dropped in the landed declaration:
    private theorem matsumura_isRegular_of_linearIndependent_cotangent
        {A : Type u} [CommRing A] [IsLocalRing A] [IsNoetherianRing A]
        [IsRegularLocalRing A]
        (n : ℕ) (rs : Fin n → A)
        (hrs_mem : ∀ i, rs i ∈ IsLocalRing.maximalIdeal A)
        (h_lin : LinearIndependent (IsLocalRing.ResidueField A)
                   (fun i => (IsLocalRing.maximalIdeal A).toCotangent
                               ⟨rs i, hrs_mem i⟩)) :
        RingTheory.Sequence.IsRegular A (List.ofFn rs) := …
    \end{verbatim}

    \emph{Proof structure} (~30--50 LOC, induction on \texttt{Fin n}
    via \texttt{Fin.induction} or \texttt{Nat} on \texttt{n}, with
    \(A\) generalized):
    \begin{enumerate}
      \item \textbf{Base case} \(n = 0\): \texttt{List.ofFn} on an
        empty \texttt{Fin 0} is \texttt{[]}, and the empty list is
        \texttt{IsRegular} trivially via \texttt{RingTheory.Sequence.isRegular\_nil}.
      \item \textbf{Inductive step} \(n = k+1\): unpack
        \(\texttt{rs}\,0\) (the head) and the tail via
        \texttt{Fin.tail}/\texttt{Fin.cases}. Apply
        \(\texttt{RingTheory.Sequence.isRegular\_cons\_iff}\)
        (\texttt{Mathlib/RingTheory/Regular/RegularSequence.lean:246}):
        the cons is \texttt{IsRegular} iff (a) \texttt{IsSMulRegular
        A (rs 0)} and (b) \texttt{IsRegular} on the tail descended
        to \(A / (rs\,0)\) on the appropriate image module.
      \item \textbf{Discharge (a)} \texttt{IsSMulRegular A (rs 0)}:
        in an \texttt{IsDomain}, \texttt{IsSMulRegular} on a non-zero
        element is automatic via \texttt{IsSMulRegular.of\_ne\_zero}
        (or its analogous Mathlib lemma — verify the name).
        Domain witness: invoke
        \(\texttt{RingTheory.CohenMacaulay.isDomain\_of\_regularLocal A}\)
        (newly-public iter-202 cross-file import). Non-zero: from
        \(\texttt{rs}\,0 \notin \mathfrak m^2\) (consequence of
        \texttt{h\_lin} on the cotangent class), \(\texttt{rs}\,0
        \neq 0\) (the zero element lies in every power of
        \(\mathfrak m\)).
      \item \textbf{Apply the RLR-preservation helper} (newly-public
        iter-202): \(\texttt{RingTheory.CohenMacaulay.regularLocal\_quotient\_isRegularLocal\_of\_notMemSq A hdim (rs 0) hxMem hxNotSq}\)
        returns a \texttt{Nontrivial} witness, \texttt{IsLocalRing}
        witness, \texttt{IsRegularLocalRing} witness, and
        \texttt{spanFinrank = k} for \(A / (rs\,0)\). Promote these
        to typeclass instances via \texttt{haveI}.
      \item \textbf{Dimension descent}: by Mathlib's
        \texttt{ringKrullDim\_quotient\_span\_singleton\_succ\_eq\_ringKrullDim}
        (\texttt{Mathlib/RingTheory/KrullDimension/Regular.lean}),
        \(\texttt{ringKrullDim}\,(A / (rs\,0)) + 1 = \texttt{ringKrullDim}\,A = n + 1\),
        so \(\texttt{ringKrullDim}\,(A / (rs\,0)) = k\).
        (Equivalently from the helper's \texttt{spanFinrank = k} +
        Mathlib's RLR Krull-dim formula.)
      \item \textbf{Apply IH on the quotient}: at the smaller
        \(n = k\), \texttt{IsRegular} on the tail image module in
        \(A / (rs\,0)\).
      \item \textbf{Close}: combine (3) + (6) into the goal via
        \texttt{isRegular\_cons\_iff.mpr}.
    \end{enumerate}

    \emph{Linear-independence transfer}: the \texttt{h\_lin}
    hypothesis on \texttt{rs} in the IH for \(n = k\) needs to be
    derived from the same \texttt{h\_lin} on \(A\) via the image of
    cotangent classes through the quotient algebra map; this is a
    straightforward \(\texttt{LinearIndependent.image}\) /
    \(\texttt{LinearIndependent.map}\) application using the
    helper's \texttt{spanFinrank} witness.

    \emph{Target sorry-line for closure}: this declaration is a NEW
    \texttt{private theorem}; its proof body has no sorry-line
    target. The iter-203 A3 consumer (existing
    \texttt{private theorem Algebra.SubmersivePresentation.relations\_isRegular\_in\_localization}
    if landed, OR composed inline in the L1061 closure of
    \texttt{isRegularLocalRing\_stalk\_of\_smooth}) invokes A1 with
    the lin-indep premise from A2 to produce the
    \texttt{RingTheory.Sequence.IsRegular} witness consumed by
    \texttt{ringKrullDim\_quotient\_localization\_MvPolynomial\_of\_regular}
    (iter-201 L924).

    \emph{Verification}: after landing, run \texttt{lean\_verify} on
    \texttt{matsumura\_isRegular\_of\_linearIndependent\_cotangent}
    and confirm \texttt{\#print axioms} returns the kernel triple
    \(\{\texttt{propext}, \texttt{Classical.choice}, \texttt{Quot.sound}\}\).
  \item \textbf{A2 — Cotangent-localisation transport}
    (sub-pieces partially landed iter-201; residual
    \(\sim 50\)--\(100\) LOC): transport
    \(\texttt{Algebra.SubmersivePresentation.basisCotangent}\)'s
    linear-independence statement from \(I/I^2\) over \(S\) to
    \(\mathfrak m_A / \mathfrak m_A^2\) over \(\kappa(\mathfrak m_A)\)
    on the localisation \(A = P.\texttt{Ring}_{m'}\). Uses
    \(\texttt{LinearIndependent.of\_isLocalizedModule\_of\_isRegular}\)
    (in \texttt{Mathlib/RingTheory/Localization/Module.lean}) +
    the conormal-localisation iso
    \((I/I^2) \otimes_S S_{\mathfrak m} \simeq
     (I \cdot A) / (I \cdot A)^2\).

    \textbf{Iter-201 landed substrate.} Two A2 sub-pieces are
    axiom-clean private theorems in
    \texttt{CodimOneExtension.lean}:
    \begin{itemize}
      \item \(\texttt{submersivePresentation\_relation\_cotangent\_mk\_linearIndependent}\)
        (L854): forward-ergonomic re-package of
        \(\texttt{Basis.linearIndependent}\) +
        \(\texttt{basisCotangent\_apply}\), expressing lin-indep
        directly in the \(\texttt{Cotangent.mk}\)-of-relation form.
      \item \(\texttt{submersivePresentation\_relation\_cotangent\_mk\_linearIndependent\_localized}\)
        (L889): transport through the canonical localisation map
        \(\texttt{LocalizedModule.mkLinearMap}\) at any prime
        \(p\) of \(S\), via
        \(\texttt{LinearIndependent.of\_isLocalizedModule}\).
    \end{itemize}
    \textbf{Iter-202+ A2 residual.} The conormal-localisation iso
    for \(\texttt{IsLocalization.AtPrime}\) (Stacks 02JK at a
    non-Away prime; the Mathlib infrastructure
    \(\texttt{Algebra.Generators.cotangentCompLocalizationAwayEquiv}\)
    generalises only to the Away case) is the remaining substantive
    gap on the A2 side. Estimated \(\sim 50\)--\(100\) LOC.
  \item \textbf{A3 — Assembly} (\(\sim 10\)--\(15\) LOC): invoke
    A1 with the lin-indep premise from A2, then close via
    \(\texttt{IsWeaklyRegular.isRegular\_of\_isLocalizedModule\_of\_mem}\)
    (in \texttt{Mathlib/RingTheory/Regular/Flat.lean:65--73} ---
    DIRECT INVOCATION). Final declaration:
    \(\texttt{Algebra.SubmersivePresentation.relations\_isRegular\_in\_localization}\).

    \textbf{Iter-201 landed substrate.}
    \(\texttt{ringKrullDim\_quotient\_localization\_MvPolynomial\_of\_regular}\)
    (\texttt{CodimOneExtension.lean} L924) is the
    \texttt{IsRegular}-conditional Steps 1+2+3 composite for the
    polynomial-ring setting: given an \(\texttt{IsRegular}\,A\,rs\)
    witness (the iter-202+ A1/A2/A3 endpoint), this lemma reduces
    the Stage 6.A conclusion to a one-line invocation.
\end{itemize}
consuming Mathlib's existing
\(\texttt{Algebra.SubmersivePresentation.jacobian\_isUnit}\) +
\(\texttt{Algebra.SubmersivePresentation.basisCotangent}\) +
\(\texttt{Algebra.SubmersivePresentation.cotangentEquiv}\) +
\(\texttt{IsWeaklyRegular.isRegular\_of\_isLocalizedModule\_of\_mem}\) +
\(\texttt{LinearIndependent.of\_isLocalizedModule\_of\_isRegular}\) +
\(\texttt{IsLocalRing.isRegular\_iff\_isWeaklyRegular\_of\_subset\_maximalIdeal}\);
(iv) plug into \texttt{ringKrullDim\_quotient\_add\_eq\_of\_regular\_sequence}
+ \texttt{ringKrullDim\_localization\_atMaximal\_MvPolynomial} for
\(\texttt{ringKrullDim}\,S_{m} = n\) (with the standard-smooth
presentation's \texttt{relativeDimension} witness as \(n\)). Estimated
project-side cost: \(\sim 30-60\) LOC for (ii)+(iii); the full L1061
inline closure additionally requires scheme-to-algebra bridging
(extracting the presentation from
\(\Gamma(X.\texttt{left}, V)\) over \(\Gamma(\Spec, U) = \bar k\),
identifying the maximal ideal of \(\Gamma(X.\texttt{left}, V)\)
corresponding to \(z\)) which adds another \(\sim 30-50\) LOC. The
substrate landing trigger Lane T32 re-engagement (per
\texttt{PROGRESS.md} Lane T32 binding trigger).

\subsection{Substrate lemmas for Stages 00OE / 00SW and the scheme bridges}
\label{subsec:codim1_substrate_lemmas}

The narrative above is realised on the Lean side by the named substrate lemmas
collected here. They feed the Stacks-00OE Krull-dimension formula
\cref{lem:smooth_algebra_krull_dim_formula} and the regular-stalk assembly,
either directly or through the Jacobian regular-sequence witness.

\paragraph{Stacks 00OE Krull-dimension chain.}

\begin{lemma}\leanok
[00OE Step 1 --- localised Krull dimension equals height]
  \label{lem:ringKrullDim_localization_eq_height_atPrime}
  \lean{AlgebraicGeometry.Scheme.ringKrullDim_localization_eq_height_atPrime}
  For a prime ideal \(\mathfrak m\) of a commutative ring \(R\) and any localisation
  \(A\) of \(R\) at \(\mathfrak m\), the Krull dimension of \(A\) equals the height of
  \(\mathfrak m\).
\end{lemma}

\begin{proof}
  Proved directly in Lean: re-export of Mathlib's
  \texttt{IsLocalization.AtPrime.ringKrullDim\_eq\_height}. A sub-step of
  \cref{lem:ringKrullDim_localization_atMaximal_mvPolynomial}.
\end{proof}

\begin{lemma}\leanok
[00OE Step 2 (lower bound) --- height $\ge n$ in $k[\mathrm{Fin}\,n]$]
  \label{lem:mvPolynomial_maximalIdeal_height_ge_card}
  \lean{AlgebraicGeometry.Scheme.MvPolynomial.maximalIdeal_height_ge_card_of_field}
  For a field \(k\), every maximal ideal of the polynomial ring
  \(\mathrm{MvPolynomial}\,(\mathrm{Fin}\,n)\,k\) has height at least \(n\).
\end{lemma}

\begin{proof}
  Proved directly in Lean: induction on \(n\), transporting the ideal along the
  one-variable-splitting isomorphism and using that adjoining a polynomial variable
  raises height by one, together with the Jacobson contraction of the maximal ideal.
  A sub-step of \cref{lem:mvPolynomial_maximalIdeal_height_eq_card}.
\end{proof}

\begin{lemma}[00OE Step 2 (upper bound) --- height $\le \#\iota$ in $k[\iota]$]
  \label{lem:mvPolynomial_maximalIdeal_height_le_natCard}
  \lean{AlgebraicGeometry.Scheme.MvPolynomial.maximalIdeal_height_le_natCard_of_field}
  For a field \(k\) and a finite index set \(\iota\), every proper ideal of
  \(\mathrm{MvPolynomial}\,\iota\,k\) has height at most \(\#\iota\).
\end{lemma}

\begin{proof}
  Proved directly in Lean: height is bounded by the Krull dimension, which for the
  Noetherian polynomial ring over a field equals \(\#\iota\). A sub-step of
  \cref{lem:mvPolynomial_maximalIdeal_height_eq_card}.
\end{proof}

\begin{lemma}\leanok
[00OE Step 2 ($\mathrm{Fin}\,n$ form) --- maximal-ideal height equals $n$]
  \label{lem:mvPolynomial_maximalIdeal_height_eq_card}
  \lean{AlgebraicGeometry.Scheme.MvPolynomial.maximalIdeal_height_eq_card}
  \uses{lem:mvPolynomial_maximalIdeal_height_ge_card, lem:mvPolynomial_maximalIdeal_height_le_natCard}
  For a field \(k\), every maximal ideal of
  \(\mathrm{MvPolynomial}\,(\mathrm{Fin}\,n)\,k\) has height exactly \(n\).
\end{lemma}

\begin{proof}
  Proved directly in Lean: combine the lower bound
  \cref{lem:mvPolynomial_maximalIdeal_height_ge_card} and the upper bound
  \cref{lem:mvPolynomial_maximalIdeal_height_le_natCard} by antisymmetry.
\end{proof}

\begin{lemma}[00OE Step 2 (general finite form) --- maximal-ideal height equals $\#\iota$]
  \label{lem:mvPolynomial_maximalIdeal_height_eq_natCard}
  \lean{AlgebraicGeometry.Scheme.MvPolynomial.maximalIdeal_height_eq_natCard}
  \uses{lem:mvPolynomial_maximalIdeal_height_eq_card}
  For a field \(k\) and a finite index set \(\iota\), every maximal ideal of
  \(\mathrm{MvPolynomial}\,\iota\,k\) has height exactly \(\#\iota\).
\end{lemma}

\begin{proof}
  Proved directly in Lean: transport
  \cref{lem:mvPolynomial_maximalIdeal_height_eq_card} along the renaming isomorphism
  induced by a bijection \(\iota \simeq \mathrm{Fin}\,(\#\iota)\). A sub-step of
  \cref{lem:ringKrullDim_localization_atMaximal_mvPolynomial}.
\end{proof}

\begin{lemma}[00OE Steps 1+2 capstone --- localised polynomial-ring Krull dimension]
  \label{lem:ringKrullDim_localization_atMaximal_mvPolynomial}
  \lean{AlgebraicGeometry.Scheme.ringKrullDim_localization_atMaximal_MvPolynomial}
  \uses{lem:ringKrullDim_localization_eq_height_atPrime, lem:mvPolynomial_maximalIdeal_height_eq_natCard}
  For a field \(k\), a finite index set \(\iota\), and any localisation \(A\) of
  \(\mathrm{MvPolynomial}\,\iota\,k\) at a maximal ideal \(\mathfrak m\), the Krull
  dimension of \(A\) equals \(\#\iota\).
\end{lemma}

\begin{proof}
  Proved directly in Lean: compose
  \cref{lem:ringKrullDim_localization_eq_height_atPrime} (Krull dimension equals
  height) with \cref{lem:mvPolynomial_maximalIdeal_height_eq_natCard} (height equals
  \(\#\iota\)).
\end{proof}

\begin{lemma}\leanok
[00OE Step 3 (additive form) --- dimension drop by a regular sequence]
  \label{lem:ringKrullDim_quotient_add_eq_regular_sequence}
  \lean{AlgebraicGeometry.Scheme.ringKrullDim_quotient_add_eq_of_regular_sequence}
  Let \(R'\) be a Noetherian local ring with \(\mathrm{ringKrullDim}\,R' = a\) and let
  \(rs\) be a regular sequence in \(R'\) of length \(b\). Then
  \(\mathrm{ringKrullDim}\,(R'/(rs)) + b = a\). The additive form is used because the
  value semiring \(\mathrm{WithBot}\,\mathbb N_\infty\) carries no subtraction.
\end{lemma}

\begin{proof}
  Proved directly in Lean: re-export of Mathlib's
  \texttt{ringKrullDim\_add\_length\_eq\_ringKrullDim\_of\_isRegular} with the
  dimension and length values made explicit. A sub-step of
  \cref{lem:ringKrullDim_quotient_localization_mvPolynomial_regular}.
\end{proof}

\begin{lemma}[00OE Steps 1+2+3 composite --- localised polynomial-quotient dimension]
  \label{lem:ringKrullDim_quotient_localization_mvPolynomial_regular}
  \lean{AlgebraicGeometry.Scheme.ringKrullDim_quotient_localization_MvPolynomial_of_regular}
  \uses{lem:ringKrullDim_localization_atMaximal_mvPolynomial, lem:ringKrullDim_quotient_add_eq_regular_sequence}
  Let \(A\) be a Noetherian localisation of \(\mathrm{MvPolynomial}\,\iota\,k\) at a
  maximal ideal, and let \(rs\) be a regular sequence in \(A\) of length \(b\). Then
  \(\mathrm{ringKrullDim}\,(A/(rs)) + b = \#\iota\).
\end{lemma}

\begin{proof}
  Proved directly in Lean: compose the capstone
  \cref{lem:ringKrullDim_localization_atMaximal_mvPolynomial} with the
  regular-sequence dimension drop
  \cref{lem:ringKrullDim_quotient_add_eq_regular_sequence}.
\end{proof}

\paragraph{Stacks 00SW submersive-presentation cotangent independence.}

\begin{lemma}\leanok
[00SW --- relations are cotangent-independent]
  \label{lem:submersive_relation_cotangent_linearIndependent}
  \lean{AlgebraicGeometry.Scheme.submersivePresentation_relation_cotangent_mk_linearIndependent}
  For a submersive presentation \(P\) of an \(R\)-algebra \(S\), the family of
  cotangent classes of the relations \(i \mapsto [\,P.\mathrm{relation}\,i\,]\) in the
  cotangent module of \(P\) is \(S\)-linearly independent.
\end{lemma}

\begin{proof}
  Proved directly in Lean: re-packaging of the linear independence of the cotangent
  basis of \(P\). A sub-step of
  \cref{lem:submersive_relation_cotangent_linearIndependent_localized}.
\end{proof}

\begin{lemma}\leanok
[00SW --- cotangent independence transported to a localisation]
  \label{lem:submersive_relation_cotangent_linearIndependent_localized}
  \lean{AlgebraicGeometry.Scheme.submersivePresentation_relation_cotangent_mk_linearIndependent_localized}
  \uses{lem:submersive_relation_cotangent_linearIndependent}
  With \(P\) a submersive presentation of \(S\) and \(p\) a prime ideal of \(S\), the
  image of the cotangent classes of the relations under the canonical localisation map
  of the cotangent module at \(p\) is \(\mathrm{Localization.AtPrime}\,p\)-linearly
  independent.
\end{lemma}

\begin{proof}
  Proved directly in Lean: transport
  \cref{lem:submersive_relation_cotangent_linearIndependent} through the localisation
  map, using that linear independence is preserved by a localised-module map.
\end{proof}

\paragraph{Relative dimension and scheme-to-algebra bridges.}

\begin{lemma}\leanok
[Stage 6 sub-gap (i) --- relative dimension of a standard-smooth algebra]
  \label{lem:exists_standardSmoothOfRelativeDimension_of_standardSmooth}
  \lean{AlgebraicGeometry.Scheme.exists_isStandardSmoothOfRelativeDimension_of_isStandardSmooth}
  Every standard-smooth \(R\)-algebra \(S\) is standard smooth of some specific
  relative dimension \(n \in \mathbb N\).
\end{lemma}

\begin{proof}
  Proved directly in Lean: read the dimension off the underlying submersive
  presentation. A sub-step of \cref{lem:smooth_to_regular_local_ring}.
\end{proof}

\begin{lemma}\leanok
[Step B.a --- submersive presentation of a relative-dimension-$n$ algebra]
  \label{lem:exists_submersivePresentation_of_standardSmoothReldim}
  \lean{AlgebraicGeometry.Scheme.exists_submersivePresentation_of_isStandardSmoothOfRelativeDimension}
  \uses{lem:exists_standardSmoothOfRelativeDimension_of_standardSmooth}
  An \(R\)-algebra \(S\) that is standard smooth of relative dimension \(n\) admits an
  explicit submersive presentation \(P\) with \(P.\mathrm{dimension} = n\).
\end{lemma}

\begin{proof}
  Proved directly in Lean: re-export of the underlying-presentation field of the
  relative-dimension instance.
\end{proof}

\begin{lemma}\leanok
[Step B.b --- affine-open stalk is a localisation]
  \label{lem:isLocalization_atPrime_stalk_of_affineOpen}
  \lean{AlgebraicGeometry.Scheme.isLocalization_atPrime_stalk_of_affineOpen}
  For an affine open \(V \ni z\) of a scheme \(X\), the stalk \(\mathcal O_{X,z}\) is
  the localisation of the section ring \(\Gamma(X, V)\) at the prime ideal of \(z\).
\end{lemma}

\begin{proof}
  Proved directly in Lean: re-export of \texttt{IsAffineOpen.isLocalization\_stalk}.
\end{proof}

\begin{lemma}\leanok
[Step B.c helper --- nonempty open on a one-point space is everything]
  \label{lem:open_eq_top_of_subsingleton}
  \lean{AlgebraicGeometry.Scheme.open_eq_top_of_subsingleton}
  On a scheme whose underlying space is a subsingleton (at most one point), every
  nonempty open subset equals the whole space.
\end{lemma}

\begin{proof}
  Proved directly in Lean: the unique-point argument. A sub-step of
  \cref{lem:gammaSpecField_ringEquiv}.
\end{proof}

\begin{lemma}\leanok
[Step B.c --- sections of $\Spec \bar k$ over a nonempty open are $\bar k$]
  \label{lem:gammaSpecField_ringEquiv}
  \lean{AlgebraicGeometry.Scheme.gammaSpecField_ringEquiv}
  \uses{lem:open_eq_top_of_subsingleton}
  For a field \(\bar k\) and any nonempty open \(U\) of \(\Spec(\bar k)\), the section
  ring \(\Gamma(\Spec \bar k, U)\) is ring-isomorphic to \(\bar k\). This gives the
  algebraically-closed base ring of the standard-smooth presentation a clean
  \(\bar k\)-identification.
\end{lemma}

\begin{proof}
  Proved directly in Lean: \(\Spec(\bar k)\) has a single point, so
  \cref{lem:open_eq_top_of_subsingleton} forces \(U = \top\) and the sections collapse
  to the global sections \(\Gamma(\Spec \bar k, \top) \simeq \bar k\).
\end{proof}

\paragraph{Matsumura regular-sequence bridge (Stacks 00NQ).}

\begin{lemma}\leanok
[Step A1 bridge --- module quotient equals ring quotient]
  \label{lem:quotSMulTop_quotientRing_linearEquiv}
  \lean{AlgebraicGeometry.Scheme.quotSMulTop_quotientRing_linearEquiv}
  For \(r\) in a commutative ring \(A\), there is a canonical \((A/(r))\)-linear
  isomorphism between the module quotient \(A/(r \cdot A)\) and the ring quotient
  \(A/(r)\).
\end{lemma}

\begin{proof}
  Proved directly in Lean: compose the tensor description of the module quotient with
  the right-unitor and promote scalars along the surjective quotient map. A sub-step of
  \cref{lem:isRegular_cons_of_quotient_ring}.
\end{proof}

\begin{lemma}\leanok
[Step A1 bridge --- ring-quotient cons rule for regular sequences]
  \label{lem:isRegular_cons_of_quotient_ring}
  \lean{AlgebraicGeometry.Scheme.isRegular_cons_of_quotient_ring}
  \uses{lem:quotSMulTop_quotientRing_linearEquiv}
  Let \(r \in A\) be such that multiplication by \(r\) is injective on \(A\), and
  suppose the images of a list \(rs\) form a regular sequence over the ring quotient
  \(A/(r)\). Then \(r :: rs\) is a regular sequence over \(A\).
\end{lemma}

\begin{proof}
  Proved directly in Lean: the module-quotient cons rule transported across the
  identification \cref{lem:quotSMulTop_quotientRing_linearEquiv}. A sub-step of
  \cref{lem:matsumura_isRegular_of_linearIndependent_cotangent}.
\end{proof}

\begin{lemma}\leanok
[Step A1 --- cotangent independence descends to the quotient]
  \label{lem:matsumura_descent_cotangent}
  \lean{AlgebraicGeometry.Scheme.matsumura_descent_cotangent}
  Let \((A, \mathfrak m)\) be a local ring and \(f_0,\dots,f_m \in \mathfrak m\) with
  \(\kappa(A)\)-linearly independent cotangent classes. Then the tail
  \(f_1,\dots,f_m\), pushed into the quotient \(A' = A/(f_0)\), has
  \(\kappa(A')\)-linearly independent cotangent classes in
  \(\mathfrak m'/\mathfrak m'^2\).
\end{lemma}

\begin{proof}
  Proved directly in Lean: the cotangent map \(A \to A'\) is surjective with kernel
  spanned by the class of \(f_0\); a kernel-disjointness argument reduces a
  \(\kappa(A')\)-relation to a \(\kappa(A)\)-relation killed by the full independence
  hypothesis. A sub-step of
  \cref{lem:matsumura_isRegular_of_linearIndependent_cotangent}.
\end{proof}

\begin{lemma}\leanok
[Step A1 --- Matsumura: cotangent-independent elements form a regular sequence]
  \label{lem:matsumura_isRegular_of_linearIndependent_cotangent}
  \lean{AlgebraicGeometry.Scheme.matsumura_isRegular_of_linearIndependent_cotangent}
  \uses{lem:matsumura_descent_cotangent, lem:isRegular_cons_of_quotient_ring}
  Let \((A, \mathfrak m)\) be a Noetherian regular local ring and \(f_1,\dots,f_n \in
  \mathfrak m\) whose cotangent classes in \(\mathfrak m/\mathfrak m^2\) are
  \(\kappa(A)\)-linearly independent. Then \(f_1,\dots,f_n\) form a regular sequence in
  \(A\) (Matsumura, \emph{Commutative Ring Theory}, Thm.~14.2; Stacks 00NQ).
\end{lemma}

\begin{proof}
  Proved directly in Lean: induction on \(n\). The head \(f_1\) is a non-zero-divisor
  (a regular local ring is a domain), the cotangent independence descends to the
  quotient \(A/(f_1)\) by \cref{lem:matsumura_descent_cotangent}, and the
  ring-quotient cons rule \cref{lem:isRegular_cons_of_quotient_ring} reassembles the
  regular sequence. In the standard-smooth application the linear-independence input is
  supplied by the relations of a submersive presentation via
  \cref{lem:submersive_relation_cotangent_linearIndependent_localized}.
\end{proof}

\begin{lemma}


\leanok
  [Smooth over \(\bar k\) \(\Rightarrow\) regular local stalk (Stacks 00TT)]
  \label{lem:smooth_to_regular_local_ring}
  \lean{AlgebraicGeometry.Scheme.isRegularLocalRing_stalk_of_smooth}
  \uses{lem:module_free_kaehler_localization, lem:cotangent_kahler_over_field,
        lem:smooth_algebra_krull_dim_formula, lem:stage6_regular_stalk_assembly,
        lem:stalkMap_flat_of_smooth, lem:exists_algebra_standardSmooth_stalk_localization,
        lem:finrank_cotangentSpace_of_bijective_residue, lem:gammaSpecField_ringEquiv,
        lem:exists_standardSmoothOfRelativeDimension_of_standardSmooth}
  % SOURCE: [Stacks Project], tag 00TT --- Jacobian criterion
  % (smooth-implies-regular direction) (read from references/stacks-algebra.tex,
  % L38593--38611; tag verified against the Stacks website tag/00TT page,
  % "Lemma 10.140.3" in section 10.140 "Smooth algebras over fields").
  % SOURCE QUOTE: "Let $k$ be any field. Let $S$ be a finite type $k$-algebra.
  % Let $X = \Spec(S)$. Let $\mathfrak q \subset S$ be a prime corresponding to
  % $x \in X$. The following are equivalent:
  % \begin{enumerate}
  % \item The $k$-algebra $S$ is smooth at $\mathfrak q$ over $k$.
  % \item We have
  % $\dim_{\kappa(\mathfrak q)} \Omega_{S/k} \otimes_S \kappa(\mathfrak q)
  % \leq \dim_x X$.
  % \item We have
  % $\dim_{\kappa(\mathfrak q)} \Omega_{S/k} \otimes_S \kappa(\mathfrak q)
  % = \dim_x X$.
  % \end{enumerate}
  % Moreover, in this case the local ring $S_{\mathfrak q}$ is regular."
  % NOTE (iter-186 plan): the iter-185 writer left an empty
  % uses-annotation here intentionally to signal "no project-side label
  % prerequisites" — the Stacks 00TT prerequisites (smooth ring map,
  % regular local ring) live in Mathlib as typeclasses
  % (Algebra.Smooth, IsRegularLocalRing) but are not packaged as
  % blueprint-labeled definitions in this project; the project's
  % def:codim_one_indeterminacy and Auslander--Buchsbaum-side
  % cor:regular_cohen_macaulay are adjacent context, not strict
  % prerequisites. blueprint-doctor (iter-185) flagged the empty
  % argument as a CRITICAL plastex/depgraph build-breaker and removed
  % it; the iter-186 plan agent further reworded this NOTE because the
  % literal backtick-quoted token was still tripping the doctor regex.
  % NOTE (iter-194 reviewer): the Lean body of
  % isRegularLocalRing_stalk_of_smooth is currently a typed sorry
  % gated on the Stacks 00OE (smooth-algebra dimension formula) +
  % Stacks 02JK (cotangent / Kahler over a field) Mathlib gaps. The
  % iter-193 prover landed Stages 5a and 5b axiom-clean (the named
  % private helpers module_free_kaehlerDifferential_localization and
  % rank_kaehlerDifferential_localization_eq_relativeDimension,
  % blueprint-pinned at \cref{lem:module_free_kaehler_localization}
  % and \cref{lem:rank_kaehler_localization_eq_relative_dim} below),
  % which set up the Kahler-differentials chain.
  % NOTE (iter-198 writer): the remaining Stage 6 gap is now
  % decomposed into three named sub-lemmas in
  % \cref{subsec:stage6_subgap_decomposition} above:
  % (6.A) \cref{lem:smooth_algebra_krull_dim_formula} (Stacks 00OE
  %       smooth-algebra Krull-dim formula),
  % (6.B) \cref{lem:cotangent_kahler_over_field} (Stacks 02JK cotangent
  %       / Kahler over a field), and
  % (6.C) \cref{lem:stage6_regular_stalk_assembly} (assembly to
  %       regular local stalk).
  % Each sub-lemma carries its own intended \lean{...} pin name and
  % Mathlib-API audit; landing 6.A--6.C axiom-clean automatically
  % closes the L526 sorry inside the Lean helper
  % isRegularLocalRing_stalk_of_smooth and propagates through the
  % downstream cascade documented in
  % \cref{subsec:stage6_subgap_decomposition}'s "Cascade to consumers"
  % paragraph.
  \textit{Source: [Stacks Project], tag 00TT.}
  Let \(X\) be a variety over an algebraically closed field \(\bar k\) (geometrically
  integral, separated, locally of finite type) and assume the structure morphism
  \(X \to \Spec(\bar k)\) is \emph{smooth}. Then for every point \(z \in X\) the
  stalk \(\mathcal O_{X, z}\) is a regular local ring.

  The hypothesis \([IsAlgClosed\,\bar k]\) is load-bearing: tag 00TT is stated
  over an arbitrary field, but the variant Mathlib will be able to discharge
  axiom-clean at commit \texttt{b80f227} routes through residue fields that are
  \(\bar k\)-rational (every closed point of \(X\) has residue field \(\bar k\) under
  algebraic closure, and the chain of localisations from a closed point reaches
  every \(z \in X\) by the standard chain-of-primes argument). Without algebraic
  closure one needs the geometric-regularity refinement of tag 00TT (Stacks tag
  \texttt{056T} in the project's \texttt{stacks-varieties.md} card) and an extra
  base-change-to-\(\bar k\) step.
\end{lemma}

\begin{proof}
  \textit{Source: [Stacks Project], tag 00TT.} The cited lemma packages the
  Jacobian-criterion equivalence and concludes ``in this case the local ring
  \(S_{\mathfrak q}\) is regular''. Translated to the scheme level: a smooth
  finite-type morphism \(X \to \Spec(\bar k)\) restricted to an affine chart
  \(\Spec(S) \subseteq X\) gives an affine algebra \(\bar k \to S\) that is
  smooth at every prime \(\mathfrak q \subset S\); tag 00TT then yields
  regularity of the localisation \(S_{\mathfrak q} = \mathcal O_{X, z}\) at every
  point \(z \in X\) lying over \(\mathfrak q\).

  Mathlib hooks the prover will assemble for the formalised version of this
  block at commit \texttt{b80f227}: (i) \texttt{Algebra.FormallySmooth} and the
  smooth-locus-of-a-scheme typeclass instance threading
  \texttt{[Smooth\,X.hom]} down to the affine algebra
  \(\bar k \to \mathcal O_X(U)\) on each open affine \(U \subseteq X\); (ii) the
  Jacobian-criterion refinement of \texttt{Module.FormallyEtale} /
  \texttt{Algebra.Smooth} via the cotangent complex \texttt{H\^{}1\,LSm} (or
  equivalent), as encoded in Mathlib's \texttt{Algebra.Smooth.iff} family; (iii)
  the closed-point selection that, given a non-closed \(z\), picks a closed
  specialisation \(x \rightsquigarrow z\) whose stalk localises to
  \(\mathcal O_{X, z}\) (Mathlib \texttt{Scheme.localRing\_at\_closed\_point}
  family); (iv) preservation of regularity under localisation of a regular
  local ring at a prime (Mathlib
  \texttt{IsRegularLocalRing.localization\_isRegularLocalRing}, or the same
  fact phrased via Stacks tag \texttt{00OF}). The composition of (i)--(iv) is
  the formalised content of the smooth-implies-regular implication; the only
  Mathlib gap at \texttt{b80f227} is in step (ii)--(iii), specifically the
  Jacobian-criterion scheme-level packaging that promotes
  \texttt{[Smooth\,X.hom]\,[IsAlgClosed\,\bar k]} to
  \texttt{IsRegularLocalRing\,(X.left.presheaf.stalk\,z)}.

  As a backup citation, Matsumura, \emph{Commutative Ring Theory}, Chapter~19
  (``Regular Rings'') develops the regular-local-ring theory needed for
  step~(iv); see \cref{chap:Albanese_AuslanderBuchsbaum} for the depth-/CM-side
  references the project already maintains.
\end{proof}

\begin{lemma}\leanok
[Smooth at a codim-1 point $\Rightarrow$ principal nonzero maximal ideal]
  \label{lem:smooth_codim_one_maximalIdeal_principal}
  \lean{AlgebraicGeometry.Scheme.smooth_codim_one_maximalIdeal_isPrincipal_and_ne_bot}
  \uses{lem:smooth_to_regular_local_ring, thm:ringKrullDim_stalk_eq_coheight}
  Let \(X\) be a smooth, geometrically-irreducible, separated, integral variety over an
  algebraically closed field \(\bar k\), and let \(z \in X\) be a point with
  \(\mathrm{coheight}\,z = 1\) (a codimension-one point). Then the maximal ideal of the
  stalk \(\mathcal O_{X,z}\) is principal and nonzero.
\end{lemma}

\begin{proof}
  Proved directly in Lean (modulo the named Stacks 00TT regularity gap of
  \cref{lem:smooth_to_regular_local_ring}): smoothness gives a regular local stalk via
  \cref{lem:smooth_to_regular_local_ring}; the coheight bridge
  \cref{thm:ringKrullDim_stalk_eq_coheight} gives Krull dimension \(1\); the
  cotangent-space characterisation of regularity then forces
  \(\dim_\kappa \mathfrak m/\mathfrak m^2 = 1\), i.e.\ a principal maximal ideal, which
  is nonzero since the stalk is not a field. This is the geometric core consumed by
  \cref{lem:smooth_codim_one_dvr}.
\end{proof}

\begin{lemma}


\leanok
  [Smooth \(\Rightarrow\) DVR at every codim-\(1\) point]
  \label{lem:smooth_codim_one_dvr}
  \lean{AlgebraicGeometry.Scheme.localRing_dvr_of_codim_one}
  \uses{lem:smooth_to_regular_local_ring, thm:ringKrullDim_stalk_eq_coheight,
        lem:smooth_codim_one_maximalIdeal_principal}
  % SOURCE: [Hartshorne], II.6, p.~130 (the definition ``regular in codimension one'',
  % together with the immediately preceding remark that on a nonsingular variety every
  % local ring is regular) (read from references/hartshorne-algebraic-geometry.pdf,
  % PDF page 147).
  % SOURCE QUOTE: "\textbf{Definition.} We say a scheme \(X\) is \emph{regular in
  % codimension one} (or sometimes \emph{nonsingular in codimension one}) if every
  % local ring \(\mathcal O_x\) of \(X\) of dimension one is regular."
  % SOURCE QUOTE: "The most important examples of such schemes are nonsingular
  % varieties over a field (I, \S5) and noetherian normal schemes. On a nonsingular
  % variety the local ring of every closed point is regular (I, 5.1), hence all the
  % local rings are regular, since they are localizations of the local rings of
  % closed points."
  % SOURCE QUOTE: "If \(Y\) is a prime divisor on \(X\), let \(\eta \in Y\) be its generic
  % point. Then the local ring \(\mathcal O_{\eta, X}\) of \(X\) at \(\eta\) is a discrete
  % valuation ring with quotient field \(K\), the function field of \(X\)."
  \textit{Source: Hartshorne, II.6, p.~130 (definition of regular in codimension one;
  the DVR consequence at a codim-1 point).}
  Let \(X\) be a nonsingular variety over \(\bar k\) and let \(\eta \in X\) be a point with
  \(\mathrm{codim}_X \overline{\{\eta\}} = 1\). Then the local ring \(\mathcal O_{X, \eta}\) is a
  discrete valuation ring with fraction field the function field \(K(X)\) of \(X\).
\end{lemma}

\begin{proof}

  % NOTE (iter-185 writer): the `` marker on this proof block reflects
  % the Krull-dim half (Tier-1 axiom-clean closure achieved iter-184 via
  % `Scheme.ringKrullDim_stalk_eq_coheight` from the iter-183
  % `Albanese/CoheightBridge.lean`, see
  % `\cref{chap:Albanese_CoheightBridge}`). The `IsRegularLocalRing` half is
  % tracked separately as the Stacks 00TT gap and is now packaged in
  % `\cref{lem:smooth_to_regular_local_ring}` above; that block carries no
  % `` because the Mathlib hooks (i)--(iv) in its proof sketch are not
  % all axiom-clean at commit `b80f227`. The composition (00TT regularity +
  % iter-183 Krull-dim bridge) inside the private Lean helper
  % `smooth_codim_one_maximalIdeal_isPrincipal_and_ne_bot`
  % (`AlgebraicJacobian/Albanese/CodimOneExtension.lean:223`) keeps a single
  % typed `sorry` on the regularity half pending an iter-185+ prover round.
  By \cref{lem:smooth_to_regular_local_ring} (Stacks tag \texttt{00TT}, the
  Jacobian-criterion direction), the stalk \(\mathcal O_{X, x}\) is regular for
  every point \(x \in X\) --- in particular at every closed point (the case
  Hartshorne~I.5.1 states classically). Localising a regular local ring at a
  prime ideal yields a regular local ring; in particular, for \(\eta \in X\) with
  \(\mathrm{codim}_X \overline{\{\eta\}} = 1\), the local ring
  \(\mathcal O_{X, \eta}\) is regular of Krull dimension \(1\). A regular local ring of Krull dimension \(1\) is exactly a
  discrete valuation ring (Atiyah--Macdonald, Proposition 9.2; or Stacks tag
  \texttt{00PD}). Its fraction field is the function field \(K(X)\) of \(X\) by
  irreducibility of \(X\) (the unique generic point of \(X\) lies above \(\eta\), and the
  stalk at the generic point of an integral scheme is the function field).

  \paragraph{Lean encoding scope.} The Lean target
  \texttt{AlgebraicGeometry.Scheme.localRing\_dvr\_of\_codim\_one} is the
  instance/lemma stating that \(\mathcal O_{X, \eta}\) is a Mathlib
  \texttt{IsDiscreteValuationRing} under the typeclass set
  \texttt{[Smooth\,(\dots)\,X.hom]\ [IsIntegral X.left]} and the hypothesis
  \texttt{Order.coheight\ \(\eta\) = 1}. The fraction-field identification with
  \texttt{X.functionField} is a downstream named instance
  \texttt{Scheme.fractionRing\_localRing\_eq\_functionField} (or whatever Mathlib name
  the prover chooses at commit \texttt{b80f227}). The regular-local-ring-of-dimension-1
  \(\Leftrightarrow\) DVR equivalence is shipped in Mathlib as
  \texttt{IsRegularLocalRing.isDiscreteValuationRing\_of\_dim\_one} (or analogous);
  the prover re-exports it rather than re-proving.
\end{proof}

\section{Milne Theorem 3.1: codimension-\(\geq 2\) extension}
\label{sec:milne_thm31}

Before the extension theorem itself we record the function-field functoriality
\(K(Y) \to K(X)\) attached to a dominant rational map, which the codim-\(1\) step of
the extension proof uses to view a representative as a morphism \(\Spec K(X) \to Y\).

\begin{lemma}\leanok
[Image of the closed point of \(\Spec K(X)\)]
  \label{lem:fromFunctionField_base_eq_genericPoint}
  \lean{AlgebraicGeometry.Scheme.RationalMap.fromFunctionField_base_eq_genericPoint}
  Let \(X\) be an integral scheme and \(Y\) an irreducible scheme, and let
  \(f \colon X \dashrightarrow Y\) be a dominant rational map. The induced morphism
  \(\Spec K(X) \to Y\) (the restriction of \(f\) to the generic point of \(X\))
  sends the closed point of \(\Spec K(X)\) to the generic point of \(Y\).
\end{lemma}

\begin{proof}\leanok
  Choose a representative \((U, \varphi_U)\) of \(f\); its underlying morphism
  \(\varphi_U \colon U \to Y\) is dominant. The morphism \(\Spec K(X) \to Y\) factors as
  \(\Spec \mathcal O_{X,\eta_X} \to U \xrightarrow{\varphi_U} Y\), where the first map
  sends the closed point to the generic point \(\eta_U\) of \(U\) (checked after the open
  immersion \(U \hookrightarrow X\), which identifies \(\eta_U\) with \(\eta_X\)). A
  dominant morphism carries the generic point of its irreducible source to a generic
  point of \(Y\); as \(Y\) is a \(T_0\) space its generic point is unique, so the image
  is \(\eta_Y\).
\end{proof}

\begin{definition}\leanok
  \label{def:functionFieldPullback}
  \lean{AlgebraicGeometry.Scheme.RationalMap.functionFieldPullback}
  \uses{lem:fromFunctionField_base_eq_genericPoint}
  For a dominant rational map \(f \colon X \dashrightarrow Y\) from an integral scheme
  \(X\) to an irreducible scheme \(Y\), the \textbf{function-field pullback} is the ring
  homomorphism \(K(Y) \to K(X)\) obtained from the germ pullback
  \(\mathcal O_{Y, f(\eta_X)} \to K(X)\) of \(f\) at the generic point, identified along
  \cref{lem:fromFunctionField_base_eq_genericPoint} with the stalk
  \(\mathcal O_{Y,\eta_Y} = K(Y)\).
\end{definition}

The difference map of Milne's Lemma~3.3 (\cref{sec:milne_lem33} below) is built by
precomposing \(f\) with the two projections \(X \times X \to X\) and pairing the results.
The precomposition operation is the following; it is the left-hand dual of the
right-composition \(\mathrm{compHom}\) already provided by the rational-map API.

\begin{definition}\leanok
  \label{def:rationalMap_precomp}
  \lean{AlgebraicGeometry.Scheme.RationalMap.precomp}
  Let \(p \colon W \to X\) be a morphism of schemes whose underlying continuous map is
  \emph{open}, and let \(f \colon X \dashrightarrow Y\) be a rational map. The
  \textbf{precomposition} \(f \circ p \colon W \dashrightarrow Y\) is the rational map
  represented, for any partial-map representative \((V, \varphi_V)\) of \(f\), by
  \[
    \bigl(p^{-1}(V),\; \; p^{-1}(V) \xrightarrow{\;p|_{p^{-1}(V)}\;} V
    \xrightarrow{\;\varphi_V\;} Y\bigr).
  \]
  The domain \(p^{-1}(V)\) is dense because the preimage of a dense open set under an open
  continuous map is dense; the same fact makes the construction independent of the chosen
  representative, so it descends to a well-defined rational map.
\end{definition}

\begin{lemma}\leanok
  \label{lem:rationalMap_precomp_compHom}
  \lean{AlgebraicGeometry.Scheme.RationalMap.precomp_compHom}
  \uses{def:rationalMap_precomp}
  Precomposition on the left commutes with right-composition by a morphism: for
  \(p \colon W \to X\) open, \(f \colon X \dashrightarrow Y\), and \(g \colon Y \to Z\),
  \[
    (f \circ p) \circ g \;=\; (f \circ g) \circ p
    \qquad\text{in}\qquad W \dashrightarrow Z .
  \]
  Furthermore, if \(p\) and \(f\) are defined over a base scheme \(S\), then so is
  \(f \circ p\).
\end{lemma}
\begin{proof}\leanok
  Both sides act on a chosen partial-map representative \((V, \varphi_V)\) of \(f\) by
  restriction along \(p\) and composition with \(g\); the resulting equality is the
  associativity of morphism composition,
  \((p|_{p^{-1}(V)} \circ \varphi_V) \circ g = p|_{p^{-1}(V)} \circ (\varphi_V \circ g)\).
  For the over-\(S\) claim, \(p|_{p^{-1}(V)}\) is an \(S\)-morphism because it commutes with
  the structure morphisms via the restriction identity
  \(p|_{p^{-1}(V)} \circ (V \hookrightarrow X) = (p^{-1}(V) \hookrightarrow W) \circ p\),
  and a composite of \(S\)-morphisms is an \(S\)-morphism.
\end{proof}

To pair the two precompositions \(f \circ \mathrm{pr}_1\) and \(f \circ \mathrm{pr}_2\) into
a single rational map \(X \times X \dashrightarrow G \times_{\bar k} G\), we use the
function-field correspondence: for an integral scheme \(X\) over \(S\) and a scheme \(Y\)
locally of finite type over \(S\), a rational map \(X \dashrightarrow Y\) over \(S\) is the
same datum as an \(S\)-morphism \(\operatorname{Spec} K(X) \to Y\) (restriction to the generic
point, spread out again by local finite type). Morphisms out of the single scheme
\(\operatorname{Spec} K(X)\) pair freely through the universal property of the fibre product.

\begin{definition}\leanok
  \label{def:rationalMap_prod}
  \lean{AlgebraicGeometry.Scheme.RationalMap.prod}
  Let \(s_X \colon X \to S\), \(s_Y \colon Y \to S\), \(s_Z \colon Z \to S\) be structure
  morphisms with \(X\) integral and \(s_Y, s_Z\) locally of finite type, and let
  \(a \colon X \dashrightarrow Y\), \(b \colon X \dashrightarrow Z\) be rational maps over
  \(S\) (i.e.\ \(a\) restricted to the generic point commutes with the structure morphisms,
  and likewise for \(b\)). The \textbf{pairing} \(a \times_S b \colon X \dashrightarrow
  Y \times_S Z\) is the rational map corresponding, under the function-field correspondence,
  to the morphism
  \[
    \operatorname{Spec} K(X)
    \xrightarrow{\;\langle a_\eta,\, b_\eta \rangle\;}
    Y \times_S Z
  \]
  obtained from the generic-point morphisms \(a_\eta \colon \operatorname{Spec} K(X) \to Y\)
  and \(b_\eta \colon \operatorname{Spec} K(X) \to Z\) by the universal property of the fibre
  product (their composites with the structure morphisms both equal
  \(\operatorname{Spec} K(X) \to \operatorname{Spec}\mathcal O_{X,\eta} \to X \to S\)).
\end{definition}

\begin{lemma}\leanok
  \label{lem:rationalMap_prod_proj}
  \lean{AlgebraicGeometry.Scheme.RationalMap.prod_compHom_fst}
  \uses{def:rationalMap_prod}
  The two projections of the pairing recover the factors, and the pairing is again over
  \(S\):
  \[
    (a \times_S b) \circ \mathrm{pr}_1 = a,
    \qquad
    (a \times_S b) \circ \mathrm{pr}_2 = b,
    \qquad
    (a \times_S b) \circ \bigl(\mathrm{pr}_1 \circ s_Y\bigr) = s_X,
  \]
  where \(\mathrm{pr}_1 \colon Y \times_S Z \to Y\), \(\mathrm{pr}_2 \colon Y \times_S Z \to Z\)
  are the projections (the last equation being the statement that \(a \times_S b\) is over
  \(S\)).
\end{lemma}
\begin{proof}\leanok
  Right-composition with a morphism is functorial on generic-point morphisms:
  \((\varphi \circ g)_\eta = \varphi_\eta \circ g\) for any rational map \(\varphi\) and
  morphism \(g\). Hence the generic-point morphism of \((a \times_S b) \circ \mathrm{pr}_1\) is
  \(\langle a_\eta, b_\eta\rangle \circ \mathrm{pr}_1 = a_\eta\), which is the generic-point
  morphism of \(a\); since a rational map out of an integral scheme is determined by its
  generic-point morphism, \((a \times_S b) \circ \mathrm{pr}_1 = a\). The second identity is
  symmetric. The over-\(S\) identity follows by combining the first identity with
  associativity of right-composition, \((a \times_S b) \circ (\mathrm{pr}_1 \circ s_Y) =
  \bigl((a \times_S b) \circ \mathrm{pr}_1\bigr) \circ s_Y = a \circ s_Y = s_X\).
\end{proof}

The first input to Milne's Theorem~3.2 is the following: a rational map from a normal
variety to a \emph{complete} variety automatically extends across every codim-\(1\)
point.

\begin{theorem}


\leanok
  [Milne Theorem~3.1: codim-\(\geq 2\) extension to a complete target]
  \label{thm:codim_one_extension}
  \lean{AlgebraicGeometry.Scheme.RationalMap.indeterminacy_codimGe2_of_smooth_of_complete}
  \uses{def:indeterminacy_locus, def:codim_one_indeterminacy,
        lem:smooth_codim_one_dvr}
  % SOURCE: [Milne, Abelian Varieties], Theorem~3.1, \S I.3, p.~16
  % (read from references/abelian-varieties.pdf, PDF page 22).
  % SOURCE QUOTE: "\textsc{Theorem 3.1.} \emph{A rational map $\varphi \colon V
  % \dashrightarrow W\( from a normal variety \)V\( to a complete variety \)W$ is defined
  % on an open subset \(U\) of \(V\) whose complement \(V \smallsetminus U\) has codimension
  % \(\geq 2\).}"
  \textit{Source: Milne, \emph{Abelian Varieties}, Theorem~3.1, \S I.3, p.~16.}
  Let \(X\) be a nonsingular variety over \(\bar k\), let \(Y\) be a complete variety over
  \(\bar k\), and let \(f \colon X \dashrightarrow Y\) be a rational map. Then the
  indeterminacy locus \(Z(f)\) (\cref{def:indeterminacy_locus}) has codimension
  \(\geq 2\) in \(X\); equivalently, \(f\) is codim-\(1\)-indeterminacy-free
  (\cref{def:codim_one_indeterminacy}).
\end{theorem}

% SOURCE QUOTE PROOF: "\textsc{Proof.} Assume first that \(V\) is a curve. Thus we are
% given a nonsingular curve \(V\) and a regular map \(\varphi \colon U \to W\) from an
% open subset of \(V\) which we want to extend to \(V\). Consider the maps
% [diagram: \(U \hookrightarrow V\), $U \xrightarrow{u \mapsto (u, \varphi(u))} Z
% \subseteq V \times W\(, \)U \xrightarrow{\varphi} W$, with projections from
% \(V \times W\) to \(V\) and to \(W\)]
% Let \(U'\) be the image of \(U\) in \(V \times W\), and let \(Z\) be its closure. The image
% of \(Z\) in \(V\) is closed (because \(W\) is complete), and contains \(U\) (the composite
% \(U \to V\) is the given inclusion), and so \(Z\) maps onto \(V\). The maps $U \to U'
% \to U\( are isomorphisms. Therefore, \)Z \to V$ is a surjective map from a complete
% curve onto a nonsingular complete curve that is an isomorphism on open subsets.
% Such a map must be an isomorphism (complete nonsingular curves are determined by
% their function fields). The restriction of the projection map \(V \times W \to W\)
% to \(Z\ (\simeq V)\) is the extension of \(\varphi\) to \(V\) we are seeking.
%   The general case can be reduced to the one-dimensional case (using schemes). Let
% \(U\) be the largest subset on which \(\varphi\) is defined, and suppose that $V
% \smallsetminus U\( has codimension 1. Then there is a prime divisor \)Z\( in \)V
% \smallsetminus U\(. Because \)V$ is normal, its associated local ring is a discrete
% valuation ring \(\mathcal O_Z\) with field of fractions \(k(V)\). The map \(\varphi\)
% defines a morphism of schemes \(\operatorname{Spec}(k(V)) \to W\), which the
% valuative criterion of properness (Hartshorne 1977, II 4.7) shows extends to a
% morphism \(\operatorname{Spec}(\mathcal O_Z) \to W\). This implies that \(\varphi\)
% has a representative defined on an open subset that meets \(Z\) in a nonempty set,
% which is a contradiction."
\begin{proof}\leanok
  Fix an indeterminacy point \(z \in Z(f)\); we must show its coheight is \(\geq 2\).
  Suppose not, so \(\mathrm{coheight}_X(z) \leq 1\), and argue by cases.

  \emph{Coheight \(0\).} Then \(z\) is a maximal point, hence \(z = \eta_X\), the generic
  point of the irreducible \(X\). But \(\eta_X\) specializes into the dense open
  \(\mathrm{Dom}(f)\), so \(z \in \mathrm{Dom}(f)\), contradicting \(z \in Z(f)\).

  \emph{Coheight \(1\).} Then \(\overline{\{z\}}\) has codimension \(1\), so by
  \cref{lem:smooth_codim_one_dvr} the local ring \(\mathcal O_{X, z}\) is a discrete
  valuation ring with fraction field the function field \(K(X)\). The rational map \(f\)
  determines a morphism \(\operatorname{Spec} K(X) \to Y\) at the generic point, and the
  over-\(\operatorname{Spec}\bar k\) condition \(f \circ Y_{\mathrm{hom}} =
  X_{\mathrm{hom}}\) makes the resulting square commute. Since \(Y\) is proper, the
  valuative criterion of properness (Milne cites Hartshorne~II.4.7) lifts
  \(\operatorname{Spec} K(X) \to Y\) across \(\operatorname{Spec} K(X) \to
  \operatorname{Spec}\mathcal O_{X, z}\) to a morphism \(\operatorname{Spec}
  \mathcal O_{X, z} \to Y\). This spreads to a regular representative of \(f\) on a
  neighbourhood of \(z\), so \(z \in \mathrm{Dom}(f)\), again a contradiction.

  Hence every point of \(Z(f)\) has coheight \(\geq 2\); equivalently \(f\) is
  codim-\(1\)-indeterminacy-free.
\end{proof}

\begin{remark}
  \label{rmk:codim1_role_of_ab}
  \textbf{Milne 3.1 is valuative-only; Auslander--Buchsbaum is not used here.} The proof
  above uses only normality (via the codim-\(1\) DVR of \cref{lem:smooth_codim_one_dvr})
  and the valuative criterion of properness (Hartshorne~II.4.7), and is
  characteristic-free. There is \emph{no} ``extend across a codim-\(\geq 2\) point'' step:
  the theorem asserts only that \(Z(f)\) has codimension \(\geq 2\), and the actual
  extension of \(f\) to a morphism happens later, once \(Z(f)\) is shown to be empty
  (\cref{thm:codim_one_extension} together with Milne~3.3 gives
  \cref{lem:av_indeterminacyLocus_eq_empty}, \(Z(f) = \emptyset\), and then
  \cref{thm:rational_map_to_av_extends} extends across the empty locus). The
  depth-\(\geq 2\)/local-cohomology (Stacks \texttt{0AVF}) route for a group-variety-free
  target is \emph{not} taken and would in any case be false for a general complete target
  (\emph{cf.}\ \(\mathbb P^2 \dashrightarrow \mathbb P^1\), whose single codim-\(2\)
  indeterminacy point admits no regular extension).

  Where \cref{cor:regular_cohen_macaulay} and the Auslander--Buchsbaum machinery of
  \cref{chap:Albanese_AuslanderBuchsbaum} \emph{do} enter is the Milne~3.3 analysis of a
  group-variety target below (regularity of the diagonal and of local quotients), not
  this codim-\(\geq 2\) statement.
\end{remark}

\section{Milne Lemma 3.3: indeterminacy locus of a map to a group variety is codim 1}
\label{sec:milne_lem33}

The second input to Milne's Theorem~3.2 is structural: when the target \(Y\) is a group
variety, the indeterminacy locus is either empty or of \emph{pure} codimension \(1\).
Combined with \cref{thm:codim_one_extension}'s ``codimension \(\geq 2\)'' conclusion,
this forces the locus to be empty.

\begin{lemma}


\leanok
  [Milne Lemma~3.3: pure-codim-\(1\) structure of the indeterminacy locus into a group
  variety]
  \label{lem:milne_codim1_indeterminacy}
  \lean{AlgebraicGeometry.Scheme.RationalMap.indeterminacy_pure_codim_one_into_grpScheme}
  \uses{def:indeterminacy_locus, def:codim1_cycles, def:order_at_point,
        def:principal_divisor}
  % SOURCE: [Milne, Abelian Varieties], Lemma~3.3, \S I.3, p.~17
  % (read from references/abelian-varieties.pdf, PDF page 23).
  % SOURCE QUOTE: "\textsc{Lemma 3.3.} \emph{Let \(\varphi \colon V \dashrightarrow G\)
  % be a rational map from a nonsingular variety to a group variety. Then either
  % \(\varphi\) is defined on all of \(V\) or the points where it is not defined form a
  % closed subset of pure codimension 1 in \(V\) (i.e., a finite union of prime
  % divisors).}"
  \textit{Source: Milne, \emph{Abelian Varieties}, Lemma~3.3, \S I.3, p.~17.}
  Let \(X\) be a nonsingular variety over \(\bar k\), let \(G\) be a group variety over
  \(\bar k\) (a smooth group scheme of finite type), and let \(f \colon X \dashrightarrow
  G\) be a rational map. Then either \(f\) is defined on all of \(X\), or the indeterminacy
  locus \(Z(f) \subseteq X\) is a closed subset of pure codimension \(1\) (i.e.\ a finite
  union of prime divisors).

  Equivalently: it is \emph{never} the case that \(f\) has a non-empty indeterminacy
  locus all of whose irreducible components have codimension \(\geq 2\).
\end{lemma}

% SOURCE QUOTE PROOF: "\textsc{Proof.} Define a rational map
%   \(\Phi \colon V \times V \dashrightarrow G\), $\;(x, y) \mapsto \varphi(x) \cdot
%   \varphi(y)^{-1}$.
% More precisely, if \((U, \varphi_U)\) represents \(\varphi\), then \(\Phi\) is the rational
% map represented by
%   $U \times U \xrightarrow{\varphi_U \times \varphi_U} G \times G
%   \xrightarrow{\mathrm{id} \times \mathrm{inv}} G
%   \xrightarrow{m} G$.
% Clearly \(\Phi\) is defined at a diagonal point \((x, x)\) if \(\varphi\) is defined at
% \(x\), and then \(\Phi(x, x) = e\). Conversely, if \(\Phi\) is defined at \((x, x)\), then
% it is defined on an open neighbourhood of \((x, x)\); in particular, there will be an
% open subset \(U\) of \(V\) such that \(\Phi\) is defined on \(\{x\} \times U\). After
% possible replacing \(U\) by a smaller open subset (not necessarily containing \(x\)),
% \(\varphi\) will be defined on \(U\). For \(u \in U\), the formula
%   \(\varphi(x) = \Phi(x, u) \cdot \varphi(u)\)
% defines \(\varphi\) at \(x\). Thus \(\varphi\) is defined at \(x\) if and only if \(\Phi\) is
% defined at \((x, x)\). The rational map \(\Phi\) defines a map
%   \(\Phi^\ast \colon \mathcal O_{G, e} \to k(V \times V)\).
% Since \(\Phi\) sends \((x, x)\) to \(e\) if it is defined there, it follows that \(\Phi\)
% is defined at \((x, x)\) if and only if
%   \(\operatorname{Im}(\mathcal O_{G, e}) \subset \mathcal O_{V \times V, (x, x)}\).
% Now \(V \times V\) is nonsingular, and so we have a good theory of divisors (AG,
% Chapter 12). For a nonzero rational function \(f\) on \(V \times V\), write
%   \(\operatorname{div}(f) = \operatorname{div}(f)_0 - \operatorname{div}(f)_\infty\),
% with \(\operatorname{div}(f)_0\) and \(\operatorname{div}(f)_\infty\) effective divisors
% --- note that \(\operatorname{div}(f)_\infty = \operatorname{div}(f^{-1})_0\). Then
%   $\mathcal O_{V \times V, (x, x)} = \{f \in k(V \times V) \mid
%   \operatorname{div}(f)_\infty \text{ does not contain } (x, x)\} \cup \{0\}$.
% Suppose \(\phi\) is not defined at \(x\). Then for some $f \in \operatorname{Im}
% (\varphi^\ast)\(, \)(x, x) \in \operatorname{div}(f)_\infty\(, and clearly \)\Phi$ is
% not defined at the points \((y, y) \in \Delta \cap \operatorname{div}(f)_\infty\).
% This is a subset of pure codimension one in \(\Delta\) (AG 9.2), and when we identify
% it with a subset of \(V\), it is a subset of \(V\) of codimension one passing through
% \(x\) on which \(\varphi\) is not defined."
\begin{proof}
  
  Following Milne's proof of Lemma~3.3. The argument has four sub-steps. Throughout we
  fix a representative \((U, \varphi_U)\) of \(f\).

  \emph{Sub-step 1: the difference rational map.} Define
  \[
    \Phi \colon X \times X \dashrightarrow G,
    \qquad
    \Phi(x, y) \;:=\; \varphi(x) \cdot \varphi(y)^{-1},
  \]
  via the composition
  \[
    U \times U
    \;\xrightarrow{\varphi_U \times \varphi_U}\; G \times G
    \;\xrightarrow{\mathrm{id} \times \mathrm{inv}}\; G \times G
    \;\xrightarrow{m}\; G,
  \]
  where \(\mathrm{inv} \colon G \to G\) is the group-variety inverse and
  \(m \colon G \times G \to G\) is the group law. The map \(\Phi\) is rational by
  composition: precomposing \(f\) with each projection
  \(\mathrm{pr}_i \colon X \times X \to X\) (which is open, being a base change of the
  flat, locally-finitely-presented structure morphism \(X \to \Spec \bar k\)) yields
  rational maps \(f \circ \mathrm{pr}_i \colon X \times X \dashrightarrow G\)
  (\cref{def:rationalMap_precomp}, whose over-\(\bar k\) part is
  \cref{lem:rationalMap_precomp_compHom}); pairing them into \(X \times X \dashrightarrow
  G \times G\) (\cref{def:rationalMap_prod}, whose projection and over-\(\bar k\) properties
  are \cref{lem:rationalMap_prod_proj}) and composing on the right with
  \((\mathrm{id} \times \mathrm{inv})\) then \(m\) gives \(\Phi\). Its domain is at least
  \(U \times U \subseteq X \times X\), which is dense because \(U\) is dense in \(X\).

  \emph{Sub-step 2: \(\Phi\)-defined at \((x, x)\) \(\Leftrightarrow\) \(\varphi\)-defined at
  \(x\).} If \(\varphi\) is defined at \(x\), the explicit formula \(\Phi(x, x) = e\) (the
  identity of \(G\)) shows \(\Phi\) is defined at \((x, x)\). Conversely, if \(\Phi\) is
  defined at \((x, x)\), openness of \(\mathrm{Dom}(\Phi)\) in \(X \times X\) produces an
  open set \(W \ni (x, x)\) on which \(\Phi\) is defined; by projection (taking the slice
  \(\{x\} \times W' \subseteq W\) for some open \(W' \subseteq X\), possibly shrinking
  \(W'\) so \(\varphi\) is defined on \(W'\) as well --- \(W'\) need not contain \(x\)), the
  identity
  \[
    \varphi(x) \;=\; \Phi(x, u) \cdot \varphi(u),
    \qquad u \in W',
  \]
  exhibits a regular representative of \(\varphi\) at \(x\). Hence \(x \in
  \mathrm{Dom}(\varphi)\) if and only if \((x, x) \in \mathrm{Dom}(\Phi)\).

  \emph{Sub-step 3: rephrase via pullback of local rings.} The rational map \(\Phi\)
  pulls \(\mathcal O_{G, e}\) back to a \(\bar k\)-subalgebra of the function field
  \(K(X \times X)\):
  \[
    \Phi^\ast \colon \mathcal O_{G, e} \;\longrightarrow\; K(X \times X).
  \]
  Because \(\Phi\) sends \((x, x)\) to \(e\) wherever it is defined, \(\Phi\) is defined at
  \((x, x)\) if and only if \(\Phi^\ast(\mathcal O_{G, e}) \subseteq \mathcal O_{X
  \times X, (x, x)}\) (i.e.\ every local function on \(G\) near \(e\) has no pole at
  \((x, x)\) after pullback).

  \emph{Sub-step 4: pole locus is pure codimension \(1\).} The variety \(X \times X\) is
  nonsingular (smoothness is stable under products over \(\bar k\)), so the Weil-divisor
  apparatus of \cref{def:codim1_cycles} and the order map of
  \cref{def:order_at_point} apply on \(X \times X\). For any nonzero rational function
  \(f \in K(X \times X)^\ast\) decompose its divisor as
  \[
    \mathrm{div}(f) \;=\; \mathrm{div}(f)_0 - \mathrm{div}(f)_\infty,
  \]
  with \(\mathrm{div}(f)_0\) and \(\mathrm{div}(f)_\infty\) effective and
  \(\mathrm{div}(f)_\infty = \mathrm{div}(f^{-1})_0\) (Milne's notation, matching the
  project's \(\div(f) = \sum \ord_Y(f) \cdot Y\) of
  \cref{def:principal_divisor}: \(\mathrm{div}(f)_0\) collects the prime divisors with
  \(\ord_Y(f) > 0\) and \(\mathrm{div}(f)_\infty\) those with \(\ord_Y(f) < 0\)). The local
  ring at \((x, x)\) can be characterised as
  \[
    \mathcal O_{X \times X, (x, x)}
    \;=\; \bigl\{\, f \in K(X \times X) \;\big|\;
    \mathrm{div}(f)_\infty \text{ does not contain } (x, x) \,\bigr\} \cup \{0\}.
  \]
  Suppose \(f \in K(X \times X)^\ast\) is not in \(\mathcal O_{X \times X, (x, x)}\);
  equivalently, \((x, x) \in \mathrm{supp}(\mathrm{div}(f)_\infty)\). The pole divisor
  \(\mathrm{div}(f)_\infty\) is a finite sum of prime divisors of \(X \times X\) (i.e.\ a
  closed subset of pure codimension \(1\) in \(X \times X\)). Its intersection with the
  diagonal \(\Delta_X \subseteq X \times X\) is a closed subset of pure codimension \(1\)
  in \(\Delta_X\) (this is Milne's reference to AG~9.2, the standard fact that a prime
  divisor of \(X \times X\) either contains \(\Delta_X\) or meets it in pure codimension
  \(1\); on a nonsingular variety this is an instance of Krull's principal-ideal theorem
  applied to the local equation of \(\Delta_X\) inside \(X \times X\)).

  \emph{Conclusion.} Suppose \(f \colon X \dashrightarrow G\) is not defined on all of
  \(X\). Pick \(x \in Z(f)\). By Sub-steps~2--3, there is some \(f_0 \in
  \Phi^\ast(\mathcal O_{G, e})\) with \((x, x) \in \mathrm{supp}(\mathrm{div}(f_0)_\infty)\).
  By Sub-step~4 the intersection \(\mathrm{supp}(\mathrm{div}(f_0)_\infty) \cap
  \Delta_X\) is pure codimension \(1\) in \(\Delta_X\), passing through \((x, x)\).
  Identifying \(\Delta_X\) with \(X\) via the diagonal isomorphism, the corresponding
  closed subset of \(X\) has pure codimension \(1\) and passes through \(x\). By Sub-step~2,
  every point of this codim-\(1\) subset is a point where \(\Phi\) is undefined on the
  diagonal, hence where \(\varphi\) is undefined. Thus \(Z(\varphi)\) contains a pure
  codim-\(1\) closed subset passing through \(x\); since \(x \in Z(\varphi)\) was arbitrary,
  every irreducible component of \(Z(\varphi)\) has codimension exactly \(1\), i.e.\
  \(Z(\varphi)\) is of pure codimension \(1\).
\end{proof}

\begin{remark}
  \label{rmk:milne_lem33_charfree}
  \textbf{Characteristic-free.} The above proof uses only: (i) the group law on \(G\)
  and its inverse (formal smoothness of \(G\) as a group object); (ii) the Weil-divisor
  apparatus on the nonsingular variety \(X \times X\) (Hartshorne~II.6, or
  \cref{chap:RiemannRoch_WeilDivisor} in the project's notation); (iii) the
  decomposition of a divisor on a nonsingular variety into effective pieces, and
  Krull's principal-ideal theorem on \(X \times X\). None of these has a characteristic
  restriction, so \cref{lem:milne_codim1_indeterminacy} holds over every
  algebraically closed base field, in arbitrary characteristic. (Milne states the
  lemma in this generality at \S I.3, p.~17.)
\end{remark}

\section{Weil-divisor obstruction to codim-1 extension}
\label{sec:weil_obstruction}

The proof of \cref{thm:codim_one_extension}'s Step~1 (codim-\(1\) ruling-out) goes
through the valuative criterion. In the project's Weil-divisor language
(\cref{chap:RiemannRoch_WeilDivisor}) the same obstruction has a cleaner phrasing:
the extension fails at a prime divisor \(W\) if and only if some affine coordinate of
\(f\) around the image \(\overline{f(\eta_W)}\) has a strictly negative order
\(\ord_W(\cdot) < 0\). This is the form the Lean prover needs when consuming
\cref{def:order_at_point} and \cref{def:principal_divisor}.

\begin{theorem}

  [Weil-divisor characterisation of the codim-1 extension obstruction]
  \label{thm:weil_divisor_obstruction}
  \lean{AlgebraicGeometry.TODO.weilDivisorObstruction}
  \uses{def:codim_one_indeterminacy, def:order_at_point, lem:smooth_codim_one_dvr,
        lem:mem_domain_partial_map_reshuffle}
  % NOTE (iter-270 dag-walker): \uses now declares
  % \cref{lem:mem_domain_partial_map_reshuffle}. The equivalence (1)
  % $\Leftrightarrow$ (2) below rests on the membership-via-partial-map
  % characterisation "$\eta_W \in \mathrm{Dom}(f)$ iff some representative
  % $(U,\varphi_U)$ is defined at $\eta_W$", which is exactly that lemma
  % (also the honest Lean fallback this theorem's \lean{...} was detached
  % to in iter-179). The edge de-isolates the reshuffle in the graph.
  % NOTE (iter-179 review, lvb-checker `codimone-iter179` corrective):
  % `\lean{...}` pin DETACHED (Path D2 option (b) per the lvb-checker). The
  % substantive Hartshorne-II.6 valuative reformulation (the iff with
  % ord_W >= 0) below is the *intended* statement; the iter-179 Lane D
  % `mem_domain_iff_exists_partialMap_through_point` decl is only a 2-LOC
  % `Scheme.RationalMap.mem_domain` reshuffle that does NOT prove the
  % ord_W content, so re-pinning `\lean{...}` to it would falsely claim
  % kernel-clean formalization of the substantive Hartshorne iff. The
  % weaker reshuffle proven this iter will be pinned to its own light
  % `\begin{lemma}` block (working title `lem:mem_domain_partial_map_reshuffle`)
  % by an iter-180+ blueprint-writer dispatch; the substantive Hartshorne
  % version stays pinned here but unlinked until a future
  % `extend_iff_pullback_order_nonneg` lemma carries
  % `Scheme.RationalMap.order` machinery (pullback `RationalMap →
  % function-field ring map`, not yet in Mathlib at b80f227). See
  % task_results/lean-vs-blueprint-checker-codimone-iter179.md +
  % proof-journal/sessions/session_179/recommendations.md REC-14.
  % NOTE (iter-194 writer): no corresponding Lean declaration found in
  % CodimOneExtension.lean --- iter-195 must either rename to match an
  % existing decl or add a new Lean stub. (The iter-179 detach was
  % structural, not a \lean{...} typo: the substantive Hartshorne-II.6
  % ord_W >= 0 reformulation requires the
  % Scheme.RationalMap-to-function-field pullback machinery, which is
  % not in Mathlib at commit b80f227 and is not yet built in this
  % project. The weaker iter-179 reshuffle Lean declaration
  % mem_domain_iff_exists_partialMap_through_point is blueprint-pinned
  % at \cref{lem:mem_domain_partial_map_reshuffle} as the honest
  % fallback; thm:weil_divisor_obstruction here remains unpinned at the
  % \lean{...} level until the genuine pullback decl lands.)
  % SOURCE: [Hartshorne], II.6, pp.~130--131 (the valuation \(v_Y\) associated to a
  % prime divisor, and the characterisation of \(\mathcal O_{X, \eta_Y}\) as the
  % subring of \(K(X)\) on which \(v_Y \geq 0\)) (read from
  % references/hartshorne-algebraic-geometry.pdf, PDF pages 147--148).
  % SOURCE QUOTE: "If \(Y\) is a prime divisor on \(X\), let \(\eta \in Y\) be its generic
  % point. Then the local ring \(\mathcal O_{\eta, X}\) of \(X\) at \(\eta\) is a discrete
  % valuation ring with quotient field \(K\), the function field of \(X\). We call the
  % corresponding discrete valuation \(v_Y\) the \emph{valuation of \(Y\)}."
  % SOURCE QUOTE: "Now let \(f \in K^\ast\) be any nonzero rational function on \(X\).
  % Then \(v_Y(f)\) is an integer. If it is positive, we say \(f\) has a \emph{zero} along
  % \(Y\), of that order; if it is negative, we say \(f\) has a \emph{pole} along \(Y\), of
  % order \(-v_Y(f)\)."
  \textit{Source: Hartshorne, II.6, pp.~130--131 (valuation \(v_Y\) and the order map).}
  Let \(X\) be a nonsingular variety over \(\bar k\), let \(Y\) be a variety over \(\bar k\),
  let \(f \colon X \dashrightarrow Y\) be a rational map, and let \(W \subseteq X\) be a
  prime divisor with generic point \(\eta_W \in W\) (so
  \(\mathrm{codim}_X \overline{\{\eta_W\}} = 1\)). Let \((U, \varphi_U)\) be any representative of
  \(f\) with \(U \cap (X \smallsetminus W) \neq \emptyset\). The following are equivalent:
  \begin{enumerate}
    \item \(f\) is defined at \(\eta_W\) (i.e.\ \(\eta_W \in \mathrm{Dom}(f)\), equivalently
      \(W \not\subseteq Z(f)\)).
    \item For some (equivalently, every) affine open neighbourhood \(V \subseteq Y\) of
      the generic image of \(f\) near \(W\), every regular function \(g \in
      \mathcal O_Y(V)\) pulls back to a rational function \(\varphi_U^\ast(g) \in K(X)\)
      satisfying
      \[
        \ord_W\bigl(\varphi_U^\ast(g)\bigr) \;\geq\; 0.
      \]
  \end{enumerate}

  Consequently: \(f\) is codim-\(1\)-indeterminacy-free
  (\cref{def:codim_one_indeterminacy}) if and only if, for every prime divisor
  \(W \subseteq X\) and every affine open \(V \subseteq Y\) as above, the pullback
  \(\varphi_U^\ast(g)\) has \(\ord_W \geq 0\) for all \(g \in \mathcal O_Y(V)\).
\end{theorem}

\begin{proof}
  
  By \cref{lem:smooth_codim_one_dvr}, the local ring \(\mathcal O_{X, \eta_W}\) is a
  DVR with fraction field \(K(X)\) and uniformiser determined by \(W\); by
  \cref{def:order_at_point} its valuation is exactly \(\ord_W\). The local ring at
  \(\eta_W\) is therefore characterised inside \(K(X)\) as the subring
  \[
    \mathcal O_{X, \eta_W} \;=\; \bigl\{\, h \in K(X) \;\big|\; \ord_W(h) \geq 0
    \,\bigr\} \cup \{0\}.
  \]
  The rational map \(f\) is defined at \(\eta_W\) if and only if, for any affine open
  \(V = \operatorname{Spec} A \subseteq Y\) around the image \(\overline{f(\eta_W)}\),
  the pullback \(\varphi_U^\ast \colon A \to K(X)\) lands in
  \(\mathcal O_{X, \eta_W}\): this is the standard local characterisation of where a
  rational map is regular at a codim-\(1\) point of a smooth variety (it is the
  ``regular = no pole'' criterion). By the displayed identity, the latter is
  equivalent to \(\ord_W(\varphi_U^\ast(g)) \geq 0\) for every \(g \in A\), which is
  condition (2). The equivalence (1)~\(\Leftrightarrow\)~(2) is therefore established.

  For the consequence, \(f\) is codim-\(1\)-indeterminacy-free iff no prime divisor \(W\) is
  contained in \(Z(f)\), which by the equivalence above is iff the \(\ord_W \geq 0\)
  condition holds for every prime divisor \(W\).
\end{proof}

\begin{remark}
  \label{rmk:weil_obstruction_application}
  \textbf{Why this form is the prover-facing one.} In the consumer chapter
  \cref{chap:Albanese_Thm32RationalMapExtension} the Lean prover combines
  \cref{thm:codim_one_extension} with \cref{lem:milne_codim1_indeterminacy}; in
  isolation, the abstract codim-\(\geq 2\) statement
  \cref{thm:codim_one_extension} suffices and \cref{thm:weil_divisor_obstruction}
  is not invoked. The Weil-divisor form is the bridge to the project's
  \(\mathrm{Scheme.WeilDivisor}\) API, which is the form A.4.a's own internal sub-proofs
  use when establishing \emph{specific} extensions in genus-\(0\) classifying-map
  arguments and in the symmetric-product Albanese chapter
  \cref{chap:Albanese_AlbaneseUP}.
\end{remark}

\begin{lemma}


\leanok
  [Definedness at a prime divisor, repackaged as a through-the-point partial map]
  \label{lem:mem_domain_partial_map_reshuffle}
  \lean{AlgebraicGeometry.Scheme.RationalMap.mem_domain_iff_exists_partialMap_through_point}
  Let \(X\) and \(Y\) be varieties over \(\bar k\) and let
  \(f \colon X \dashrightarrow Y\) be a rational map. For a prime divisor
  \(W \subseteq X\) with generic point \(\eta_W = W.\mathrm{point}\), the
  following are equivalent:
  \begin{enumerate}
    \item \(\eta_W \in \mathrm{Dom}(f)\) (\(f\) is defined at the generic point of
      \(W\)).
    \item There exists a partial-map representative \((U, \varphi_U)\) of \(f\) ---
      i.e.\ a regular morphism \(\varphi_U \colon U \to Y\) on a dense open
      \(U \subseteq X\) whose induced rational map equals \(f\) --- whose domain
      \(U\) contains \(\eta_W\).
  \end{enumerate}
\end{lemma}

\begin{proof}

  The lemma is a definitional reshuffle of the Mathlib membership predicate
  \texttt{Scheme.RationalMap.mem\_domain}, which already states that
  \(\eta_W \in \mathrm{Dom}(f)\) iff there is some partial-map representative of
  \(f\) defined at \(\eta_W\); the only content here is reordering the two
  conjuncts under the existential so the \(\varphi_U \mapsto f\) component lands
  first (matching the way later witnesses in the project construct partial-map
  representatives). The forward and backward implications are immediate from
  the unfolded definition.
\end{proof}

\begin{remark}
  \label{rmk:mem_domain_partial_map_reshuffle_scope}
  \textbf{Relationship to \cref{thm:weil_divisor_obstruction}.}
  \cref{lem:mem_domain_partial_map_reshuffle} is the iter-179 Lane~D honest
  fallback for the substantive Hartshorne-II.6 ``regular = no pole''
  reformulation that \cref{thm:weil_divisor_obstruction} states in its full
  \(\ord_W \geq 0\) form. The reshuffle captures only the
  definitional half of the equivalence (existence of a partial-map
  representative through a prime-divisor generic point); upgrading it to the
  substantive \(\ord_W \geq 0\) form requires the
  \texttt{Scheme.RationalMap}-to-function-field pullback machinery, which is
  not shipped by Mathlib at commit \texttt{b80f227}. See the iter-179
  \texttt{\% NOTE} block on \cref{thm:weil_divisor_obstruction} for the
  history.
\end{remark}

\section{Lean encoding}
\label{sec:codim1_lean_encoding}

The Lean target file is the new
\texttt{AlgebraicJacobian/Albanese/CodimOneExtension.lean}. The declarations to
introduce, in order:

\begin{enumerate}
  \item \texttt{AlgebraicGeometry.Scheme.RationalMap.indeterminacyLocus} ---
    \cref{def:indeterminacy_locus}, a method on the Mathlib
    \texttt{Scheme.RationalMap} structure. Implementation: \texttt{X.carrier
    \textbackslash{} f.domain} as a \texttt{Set X}, paired with an
    \texttt{IsClosed} witness (either bundled or as a separate
    \texttt{Scheme.RationalMap.isClosed\_indeterminacyLocus} lemma).
  \item \texttt{AlgebraicGeometry.Scheme.RationalMap.CodimOneFree} ---
    \cref{def:codim_one_indeterminacy}, a \texttt{Prop}-valued predicate. Encoding:
    \texttt{\(\forall\) \(\eta\) : X, Order.coheight \(\eta\) = 1 \(\to\) \(\eta\) \(\in\)
    f.domain} (matching the \texttt{Order.coheight} idiom of
    \cref{def:prime_divisor}).
  \item \texttt{AlgebraicGeometry.Scheme.localRing\_dvr\_of\_codim\_one} ---
    \cref{lem:smooth_codim_one_dvr}, a lemma (or instance) producing
    \texttt{IsDiscreteValuationRing (Scheme.LocalRing.at\ X \(\eta\))} under the
    typeclass set \texttt{[Smooth (\(\Spec \bar k\)) X] [IsIntegral X]} and the
    hypothesis \texttt{Order.coheight \(\eta\) = 1}.
  \item \texttt{AlgebraicGeometry.Scheme.RationalMap.extend\_of\_codimOneFree\_of\_smooth}
    --- \cref{thm:codim_one_extension}, the existence-of-extension theorem. Signature
    (informal): given a smooth variety \(X\), a complete variety \(Y\), and
    \(f\)~codim-\(1\)-indeterminacy-free, produce a unique regular morphism
    \(\widetilde f \colon X \to Y\) extending \(f\). Consumes
    \cref{cor:regular_cohen_macaulay} from
    \cref{chap:Albanese_AuslanderBuchsbaum} as a Mathlib-style import.
  \item \texttt{AlgebraicGeometry.Scheme.RationalMap.indeterminacy\_pure\_codim\_one\_into\_grpScheme}
    --- \cref{lem:milne_codim1_indeterminacy}. Signature (informal): given a smooth
    variety \(X\), a group scheme \(G\) over \(\bar k\) that is a variety (smooth, finite
    type, separated), and \(f \colon X \dashrightarrow G\) rational, conclude that
    \(Z(f)\) is either empty or of pure codimension \(1\) in \(X\).
  \item % NOTE (iter-184 review): the declaration name `extend\_iff\_order\_nonneg`
        % was RETIRED in iter-179 (renamed to
        % `mem\_domain\_iff\_exists\_partialMap\_through\_point`; per the
        % iter-179 review of auditor 178B). The new declaration is at
        % AlgebraicJacobian/Albanese/CodimOneExtension.lean:492.
        % NOTE (iter-185 writer): the lightweight blueprint home for the
        % reshuffle now exists at \cref{lem:mem_domain_partial_map_reshuffle};
        % the substantive Hartshorne-II.6 `ord_W \geq 0` form remains at
        % \cref{thm:weil_divisor_obstruction} (unpinned until the
        % `Scheme.RationalMap` -> function-field pullback machinery lands).
    \texttt{AlgebraicGeometry.Scheme.RationalMap.mem\_domain\_iff\_exists\_partialMap\_through\_point}
    --- the iter-179 definitional reshuffle pinned at
    \cref{lem:mem_domain_partial_map_reshuffle}. The substantive
    Hartshorne-II.6 ``regular = no pole'' Weil-divisor reformulation lives
    at \cref{thm:weil_divisor_obstruction} (unpinned at the chapter level;
    the iter-179 detach \texttt{\% NOTE} on that theorem block explains the
    reshuffle). Re-uses
    \texttt{AlgebraicGeometry.Scheme.RationalMap.order} from
    \cref{chap:RiemannRoch_WeilDivisor} (the order map at a prime divisor).
\end{enumerate}

\paragraph{Mathlib readiness audit.} At commit \texttt{b80f227} the prover should
verify:
\begin{itemize}
  \item \texttt{Scheme.RationalMap.domain} is shipped at
    \texttt{Mathlib.AlgebraicGeometry.Birational.RationalMap}; its complement
    \texttt{indeterminacyLocus} is a one-line definition.
  \item The valuative criterion of properness in Mathlib form is
    \texttt{AlgebraicGeometry.IsProper.valuative\_criterion} (or analogous); the
    prover invokes this in Step~1 of \cref{thm:codim_one_extension}'s proof.
  \item % NOTE (iter-184 review): the previous claim that Mathlib's
        % \texttt{AlgebraicGeometry.Smooth} ``bundles the regular-local-rings
        % hypothesis'' was FALSE at b80f227 — empirically verified by
        % lean-vs-blueprint-checker iter184-codimone. The implication
        % \texttt{Smooth X.hom + [IsAlgClosed kbar] $\Rightarrow$ IsRegularLocalRing
        % (X.left.presheaf.stalk z)} is a Mathlib GAP (Stacks tag 00TT, the
        % Jacobian-criterion direction), not a bundled instance. The
        % \texttt{Krull-dim} half of \texttt{lem:smooth\_codim\_one\_dvr} now
        % closes axiom-clean via \texttt{Scheme.ringKrullDim\_stalk\_eq\_coheight}
        % (iter-183 \texttt{Albanese/CoheightBridge.lean}) + \texttt{exact\_mod\_cast};
        % the \texttt{IsRegularLocalRing} half remains as a typed sorry pending
        % Stacks 00TT formalisation (iter-185+ blueprint-expansion required).
        % Once \texttt{IsRegularLocalRing} is in scope, the closure
        % ``regular local + Krull dim 1 = DVR'' goes via
        % \texttt{IsDiscreteValuationRing.TFAE.out 0 4} (the iter-178
        % closure pattern).
        The DVR-at-codim-1 lemma reduces to ``regular local of dim 1 = DVR''
    (\texttt{IsDiscreteValuationRing.TFAE.out 0 4}; the iter-178 closure pattern).
    The bridge \texttt{Smooth X.hom + [IsAlgClosed kbar] $\Rightarrow$
    IsRegularLocalRing (stalk z)} is the Stacks 00TT gap — see the iter-184 note above.
  \item Mathlib's product-of-smooth-is-smooth instance,
    \texttt{AlgebraicGeometry.Smooth.comp\_pullback} / equivalent, is needed in
    \cref{lem:milne_codim1_indeterminacy}'s Sub-step~4 to apply the Weil-divisor
    machinery to \(X \times X\).
  \item The local-cohomology vanishing in Step~2 of
    \cref{thm:codim_one_extension}'s proof
    (\(H^0_{\{x\}}(V, \mathcal O_X) = H^1_{\{x\}}(V, \mathcal O_X) = 0\) at a
    depth-\(\geq 2\) point) is a substantial Mathlib gap as of \texttt{b80f227}; the
    prover's options are (i) prove the depth-\(\geq 2\) extension property as a local
    lemma using the \texttt{Module.depth} / Auslander--Buchsbaum machinery of
    \cref{chap:Albanese_AuslanderBuchsbaum}, or (ii) reformulate
    \cref{thm:codim_one_extension} in the weaker ``surface'' form
    (where \(X\) has \(\dim 2\), so codim-\(\geq 2\) means closed points only and the
    extension is the classical Hartshorne~III.7.1 ``smooth-surface point removal''
    argument). For the consumer \cref{chap:Albanese_Thm32RationalMapExtension} only
    the surface case is needed in practice (the symmetric product \(C^{(g)}\) is
    smooth, and the Abel--Jacobi map's intermediate-step indeterminacy lives on a
    smooth surface inside \(C^{(g)} \times C^{(g)}\)).
\end{itemize}

No new core axiom is required. The proof of \cref{lem:milne_codim1_indeterminacy} is
purely Weil-divisor manipulation on a smooth variety and is the cleanest part of
the chapter; the substance of the Lean engineering is in
\cref{thm:codim_one_extension}'s proof, both halves of which rest on Mathlib API
gaps the prover must either fill in or work around.

\section{Out of scope}
\label{sec:codim1_out_of_scope}

This chapter packages \emph{only} the two extension inputs and the Weil-divisor
reformulation needed for A.4.a. The following adjacent content is explicitly out of
scope:
\begin{itemize}
  \item \textbf{The Auslander--Buchsbaum formula and the regular \(\Rightarrow\)
    Cohen--Macaulay corollary.} Proved in \cref{chap:Albanese_AuslanderBuchsbaum}
    as \cref{thm:auslander_buchsbaum} and \cref{cor:regular_cohen_macaulay}; this
    chapter only consumes them (in Step~2 of \cref{thm:codim_one_extension}; see
    \cref{rmk:codim1_role_of_ab}).
  \item \textbf{Milne Theorem~3.2 itself.} The combination of
    \cref{thm:codim_one_extension} and \cref{lem:milne_codim1_indeterminacy} into the
    statement ``a rational map from a nonsingular variety to an abelian variety
    extends'' lives in \cref{chap:Albanese_Thm32RationalMapExtension} as
    \cref{thm:rational_map_to_av_extends}; this chapter only supplies the two inputs.
  \item \textbf{Milne Theorem~3.4 / Corollary~3.7 / Theorem~3.8 / Proposition~3.9 /
    Proposition~3.10.} The remainder of Milne~\S I.3 (the ``map between AVs is
    translated homomorphism'' family of results, and the
    \(\mathbb A^1 \dashrightarrow A\) / unirational\(\dashrightarrow A\) constancy
    statements) is downstream of \cref{thm:rational_map_to_av_extends} but is not
    needed by A.4 in the project's current routing: the rigidity corollaries are
    handled via the proven Rigidity Lemma, and the positive-genus arm consumes only
    the bare Theorem~3.2 statement.
  \item \textbf{Weil-divisor degree on \(X \times X\) for \(X\) a curve.} The proof of
    \cref{lem:milne_codim1_indeterminacy} invokes the Weil-divisor apparatus on the
    smooth surface \(X \times X\) but does \emph{not} need the degree map of
    \cref{def:divisor_degree} (which is curve-specific); only the codim-\(1\)
    decomposition \(\div(f) = \div(f)_0 - \div(f)_\infty\) on a nonsingular surface is
    used. The bilinear degree formula on \(X \times X\) for \(X\) a curve (intersection
    numbers, Hartshorne~V.1) is out of scope.
  \item \textbf{Cartier divisors and the Cartier/Weil comparison.} On a smooth variety
    they coincide and the chapter uses Weil throughout; the Cartier formalism is not
    introduced.
  \item \textbf{Generalisation to non-algebraically-closed base fields.} The
    project's positive-genus arm runs over \(\bar k\) throughout; descent of the
    extension to a non-algebraically-closed base, if ever needed, would require
    Galois-descent machinery and is not in scope.
  \item \textbf{Resolution of indeterminacy by blow-ups} (Hironaka-style, for higher
    indeterminacy codimensions on a smooth variety with a non-group-variety target):
    out of scope. The combination of \cref{thm:codim_one_extension} and
    \cref{lem:milne_codim1_indeterminacy} for \(G\) a group variety is sharper than any
    blow-up-based approach because it produces an extension on the original \(X\)
    rather than on a blow-up; no blow-up apparatus is needed.
\end{itemize}
